\documentclass[10.5pt,a4paper]{article}
\usepackage[a4paper,margin=2.2cm]{geometry}
\usepackage{fontspec}
\usepackage{graphicx,amsmath,amssymb,booktabs,xcolor,needspace,enumitem}
\usepackage[colorlinks=true,linkcolor=blue!50!black,urlcolor=blue!50!black]{hyperref}
\setlength{\parindent}{0pt}\setlength{\parskip}{5pt}
\setlist{itemsep=2pt,topsep=2pt,leftmargin=1.4em}
\newcommand{\Araw}{A_{\mathrm{raw}}}
\newcommand{\MeV}{\,\mathrm{MeV}}
\newcommand{\GeV}{\,\mathrm{GeV}}
\newcommand{\E}{\mathbb{E}}
\newcommand{\Var}{\mathrm{Var}}
\newcommand{\fig}[2]{\begin{center}\includegraphics[width=\linewidth,height=#2\textheight,keepaspectratio]{../report_complete/figures/#1.pdf}\end{center}}
\newcommand{\figh}[3]{\clearpage\needspace{0.6\textheight}\subsection*{#1}\addcontentsline{toc}{section}{#1}\fig{#2}{#3}}
\newcommand{\sez}[1]{\par\Needspace{6\baselineskip}\medskip\noindent\textbf{\large #1}\par\nopagebreak\smallskip}
\title{\textbf{Le figure del report, spiegate}\\[2pt]\large Frequentist and Bayesian Analysis of $B^\pm\to K^\pm K^+K^-$ at LHCb}
\author{Marco Belotti, Giacomo Moretti -- appunti di studio}
\date{Numerazione delle figure = quella di \texttt{report\_complete/main.pdf} (ottobre 2026)}
\begin{document}
\maketitle
\tableofcontents

\section*{Come leggere questo documento}
Ogni figura ha cinque blocchi:
\begin{enumerate}
\item \textbf{Il problema}: quale domanda dell'analisi la figura risolve e perché la domanda è necessaria.
\item \textbf{Teoria}: i concetti, le definizioni e le formule che servono per capire cosa è stato calcolato (ricavate dove sono brevi).
\item \textbf{Come si legge}: passi numerati per leggere il grafico: assi, elementi, cosa guardare, che aspetto avrebbe un risultato buono o cattivo.
\item \textbf{Cosa dice}: i risultati con i numeri del report, l'interpretazione e la conseguenza per la conclusione.
\item \textbf{Limiti e orale}: ciò che la figura non prova, i caveat, e le domande che un esaminatore può fare, con la risposta.
\end{enumerate}
I valori letti a occhio dai grafici sono indicati con $\approx$; quelli senza sono nel testo del report.

\textbf{Notazione di base.} $A\equiv\Araw=[N(B^-)-N(B^+)]/[N(B^-)+N(B^+)]$ è il parametro d'interesse (POI); gli altri undici parametri del fit sono \emph{nuisance}. Il dataset è $\{(m_i,q_i)\}_{i=1}^N$, massa invariante e carica del $B$ di ciascun candidato; $N=19\,731$ nella finestra $[5100,5600]\MeV$. Il modello nominale: segnale gaussiano (stessi $\mu,\sigma$ per $B^\pm$) più fondo a doppia esponenziale con forma propria per ciascuna carica, con
\[
N_s^\pm=\tfrac{N_s}{2}(1\mp A),\qquad N_b^\pm=\tfrac{N_b}{2}(1\mp A_b).
\]
Dodici parametri: $\theta=(N_s,N_b,A,A_b,\mu,\sigma,\lambda_1^\pm,\lambda_2^\pm,f_\pm)$.

% ======================================================================== I
\figh{Fig.\ I --- Spettro di massa invariante (Sez.\ II~B)}{mass_spectrum}{0.34}
\sez{Il problema}
Prima di qualunque fit bisogna \emph{guardare} i dati. Tre domande concrete: (i) che forma ha lo spettro, e quindi quale modello è ragionevole; (ii) le due cariche del $B$ e le due polarità del magnete hanno la stessa forma, così che si possa assumere massa e risoluzione comuni; (iii) c'è abbastanza statistica per misurare un effetto del $4\%$. La figura risponde alle prime due e quantifica la terza.

\sez{Teoria}
La massa invariante del candidato $B$ si ricostruisce dalle tre tracce, a cui si assegna l'ipotesi di massa del kaone ($M_K=493.677\MeV$) indipendentemente dal PID:
\[
m(KKK)=\sqrt{\Big(\sum_{i=1}^3E_i\Big)^2-\Big|\sum_{i=1}^3\vec p_i\Big|^2},\qquad E_i=\sqrt{|\vec p_i|^2+M_K^2}.
\]
Lo spettro è la somma di tre componenti.
\begin{itemize}
\item \emph{Segnale:} $B^\pm\to K^\pm K^+K^-$ veri. La massa ricostruita fluttua attorno a $m_B$ con una larghezza che è la \emph{risoluzione} del rivelatore (qui $\sigma\approx18\MeV$): da qui il picco quasi gaussiano. Poiché particella e antiparticella hanno la stessa massa e il rivelatore è (quasi) simmetrico, $\mu$ e $\sigma$ sono assunti uguali per $B^+$ e $B^-$.
\item \emph{Fondo combinatorio:} tre tracce non correlate danno una massa senza struttura, decrescente: un'esponenziale $E(m;\lambda)\propto e^{-\lambda m}$ troncata alla finestra.
\item \emph{Fondo parzialmente ricostruito:} decadimenti di $B$ con una particella persa (per esempio un $\pi^0$) hanno massa ricostruita \emph{inferiore} a $m_B$; producono la risalita sotto $\sim5150\MeV$ e sono la ragione del secondo termine esponenziale.
\end{itemize}
Il contenuto di ogni bin è una variabile di Poisson, con fluttuazione relativa $1/\sqrt n$. L'asimmetria di conteggio (tutti i candidati, segnale e fondo insieme) ha incertezza binomiale
\[
\sigma(A_{\rm count})=\sqrt{\frac{1-A_{\rm count}^2}{N}}=\sqrt{\frac{1-0.0432^2}{19\,731}}=0.0071.
\]

\sez{Come si legge}
\begin{enumerate}
\item Asse $x$: massa $m(KKK)$ in MeV, da $5100$ a $5600$ (la finestra di fit); asse $y$: candidati per bin da $5\MeV$.
\item Pannello (a): due istogrammi sovrapposti, $B^+$ (blu) e $B^-$ (rosso). Se le due cariche hanno la stessa forma le linee coincidono; l'asimmetria si manifesta come una \emph{piccola differenza di altezza} (qui il picco blu sta leggermente sopra il rosso).
\item Pannello (b): stessi dati divisi per polarità, magnet down (verde) e up (arancione). Le altezze sono diverse perché le statistiche lo sono; si guarda la \emph{forma}.
\item In ogni pannello cercare tre zone: il picco ($\approx5250$--$5320\MeV$), la coda piatta alta ($>5350\MeV$) e la salita a sinistra ($5100$--$5150\MeV$).
\end{enumerate}

\sez{Cosa dice}
\begin{enumerate}
\item \textbf{Forma.} Picco a $\approx5284\MeV$ di altezza $\approx900$ candidati per bin da $5\MeV$ su un fondo che scende da $\approx70$ per bin a $5100\MeV$ a $\approx7$ a $5600\MeV$: coerente con $N_b=4404$ distribuito su $100$ bin ($\approx44$ in media). Il rapporto segnale/fondo nella finestra $\pm2\sigma$ è circa $96\%$ di purezza.
\item \textbf{Cariche.} Le due curve del pannello (a) hanno la stessa forma, quindi massa e risoluzione comuni sono plausibili. Con $A_{\rm count}=-0.0432$ il rapporto $N^-/N^+=(1+A)/(1-A)\approx0.92$: un $8\%$ di candidati in meno per $B^-$, cioè $\approx70$ candidati su $900$ al picco.
\item \textbf{Perché non si vede a occhio.} In un bin la fluttuazione di Poisson è $\sqrt{900}\approx30$, il $3\%$: la differenza è $\approx2\sigma$ per bin. Sommando tutti i candidati il test binomiale esatto dà $6.1\sigma$; il fit, che usa la forma e separa il fondo, dà $5.2\sigma$ sull'asimmetria del \emph{segnale}. È la sommatoria dell'informazione su $\sim15$ bin del picco che rende la misura possibile.
\item \textbf{Polarità.} Magnet down ha $\approx1.4$ volte il campione di up ($11\,636$ contro $8095$), con la stessa forma: i due campioni sono confrontabili e le misure per polarità (Fig.~XVII) sono indipendenti.
\end{enumerate}

\sez{Limiti e orale}
\begin{itemize}
\item La figura \emph{non prova} l'asimmetria: giustifica il modello e la finestra.
\item \textbf{``Perché la finestra parte da $5100$ e non da $5000$?''} Sotto $5100\MeV$ il fondo parzialmente ricostruito cambia forma e il modello a due esponenziali non lo descrive; la finestra è variata nelle sistematiche (Tab.~II).
\item \textbf{``Perché non si usa la massa del $B$ dalla simulazione?''} Perché $\mu$ è un parametro libero del fit: il valore fittato non viene imposto.
\item Il picco non è esattamente gaussiano (Fig.~III): code più pesanti.
\end{itemize}

% ======================================================================== II
\figh{Fig.\ II --- Copertura degli intervalli di Wilks (Sez.\ IV~D, di Marco)}{frequentist_wilks_coverage}{0.30}
\sez{Il problema}
L'intervallo di confidenza frequentista al $68\%$ si ottiene tagliando il profilo di verosimiglianza a $q(A)\le1$. Quella ricetta si fonda su un teorema \emph{asintotico} (Wilks). Con $N\approx2\times10^4$ candidati e un modello a $12$ parametri vale davvero? Se non valesse, l'incertezza quotata sarebbe sbagliata. Bisogna verificarlo direttamente \emph{sui dati}.

\sez{Teoria}
\emph{Verosimiglianza estesa.} Con yield di segnale e fondo come parametri, il numero totale di candidati è Poisson con media $\nu=N_s+N_b$, e
\[
L(D\mid\theta)=\frac{e^{-\nu}\nu^N}{N!}\prod_{i=1}^N\frac{N_s^{q_i}S(m_i)+N_b^{q_i}B^{q_i}(m_i)}{\nu},
\]
cioè il prodotto di un termine di Poisson sul totale e di una verosimiglianza di forma.

\emph{Rapporto di verosimiglianza profilato.} Per il parametro d'interesse $A$ fissato a un valore $A_0$:
\[
q(A_0)=-2\ln\frac{L(A_0,\hat{\hat\theta}(A_0))}{L(\hat A,\hat\theta)},
\]
dove $\hat{\hat\theta}(A_0)$ massimizza $L$ a $A=A_0$ (il \emph{profilo}) e $(\hat A,\hat\theta)$ è il massimo globale.

\emph{Teorema di Wilks.} Sotto condizioni di regolarità (campione grande, punto interno allo spazio dei parametri, ipotesi annidate, modello identificabile), se $A_0$ è il valore vero allora $q(A_0)\to\chi^2_1$ per $N\to\infty$, \emph{qualunque} siano i valori dei nuisance. Conseguenza: $P[q(A_{\rm true})\le c]=P[\chi^2_1\le c]$; per $c=1$ vale $0.6827$ (il livello di $\pm1\sigma$), per $c=3.84$ vale $0.95$. L'insieme $\{A:q(A)\le c\}$ ha dunque copertura nominale.

\emph{Copertura} = frazione di ripetizioni dell'esperimento in cui l'intervallo contiene il valore vero. Qui le ripetizioni sono \emph{bootstrap}: si estraggono $N$ candidati \emph{con reinserimento} dal campione osservato (riproduce le fluttuazioni statistiche, anche della conta delle cariche) e si rifà il fit; il valore di riferimento delle repliche è il best-fit sui dati originali, $A_{\rm true}=-0.0441$.

\emph{Pull.} $(\hat A-A_{\rm true})/\sigma(\hat A)$ con $\sigma$ di Hesse $=\big[-\partial^2\ln L/\partial A^2\big]^{-1/2}$: media $0$ significa stimatore non polarizzato, deviazione standard $1$ significa incertezza corretta.

\sez{Come si legge}
\begin{enumerate}
\item \textbf{Pannello sinistro.} Asse $x$: $t=-2\ln(L_0/L_{\rm best})$ da $0$ a $10$; asse $y$: densità. L'istogramma azzurro è la distribuzione di $t$ valutata al valore di riferimento nelle $1000$ repliche; la curva nera tratteggiata è la $\chi^2_1$ attesa; la linea rossa punteggiata a $t=1$ è il limite al $68.3\%$.
\item L'\emph{area} dell'istogramma a sinistra della linea rossa è la \emph{copertura} dell'intervallo $q\le1$: il titolo del grafico la riporta (``Coverage = 68.4\%'').
\item Si verifica la forma: l'istogramma deve seguire la curva nera, in particolare il picco stretto vicino a $0$ (la densità $\chi^2_1\propto t^{-1/2}$ diverge) e la coda fino a $t\approx10$.
\item \textbf{Pannello destro.} Asse $x$: pull; istogramma azzurro contro la normale standard (rossa). In alto media ($-0.04\pm0.03$) e deviazione standard ($0.98$). Si verifica: centrata su $0$? larga $1$? code gaussiane?
\end{enumerate}

\sez{Cosa dice}
\begin{enumerate}
\item \textbf{Wilks vale.} L'istogramma segue la $\chi^2_1$ su tutto l'intervallo, compreso lo spike vicino a $0$.
\item \textbf{Copertura nominale.} $68.4\%$ contro $68.3\%$ atteso: scostamento di $0.1$ punti percentuali, $0.07$ errori standard (l'errore binomiale è $\sqrt{0.683\cdot0.317/1000}=1.5\%$). $0$ fit non convergiti su $1000$.
\item \textbf{Stimatore non polarizzato.} Media del pull $-0.04\pm0.03$ ($=0.98/\sqrt{1000}$), compatibile con $0$ entro $1.3$ errori standard.
\item \textbf{Incertezza corretta.} Deviazione standard $0.98$: la matrice di Hesse sovrastima l'incertezza di appena il $2\%$; se fosse $1.2$ l'incertezza quotata ($0.0085$) sarebbe sottostimata del $20\%$.
\item \textbf{Conseguenza.} L'intervallo $[-0.0526,-0.0356]$ al $68\%$ e il $5.2\sigma$ non sono artefatti di un'approssimazione asintotica a questo campione.
\end{enumerate}

\sez{Limiti e orale}
\begin{itemize}
\item Nel testo le $1000$ repliche sono bootstrap (ricampionamento dei dati), nel titolo del grafico ``toys'': vocabolario da uniformare (nota per Marco).
\item Il bootstrap ricampiona i \emph{dati}, che contengono ciò che il modello non descrive (code del segnale); i toy bayesiani della Fig.~XVIII vengono dal \emph{modello} (test di chiusura). I due controlli non sono equivalenti: questo è più severo.
\item \textbf{``Quando Wilks fallisce?''} Parametro al bordo (un yield vicino a $0$), pochi eventi, ipotesi non annidate, modello non identificabile sotto $H_0$. Qui $A\in(-1,1)$ è lontano dai bordi.
\item \textbf{``Perché non usare direttamente la Hesse?''} Perché assume che $\ln L$ sia parabolica: il profilo (Minos) non lo assume; qui coincidono (Hesse $0.0085$, Minos $[-0.0526,-0.0356]$), e il confronto è parte del controllo.
\end{itemize}

% ======================================================================== III
\figh{Fig.\ III --- Fit nominale ai due campioni di carica (Sez.\ V~C)}{bayes_nominal_mass_fit}{0.40}
\sez{Il problema}
Le incertezze su $A$ valgono solo se il modello descrive i dati. Una dichiarazione di adeguatezza richiede di vedere il confronto modello--dati \emph{carica per carica}, perché un difetto presente in una sola carica potrebbe simulare un'asimmetria.

\sez{Teoria}
Il fit è \emph{unbinned}: usa ogni $m_i$. Per confrontarlo con i dati lo si integra in bin: il valore atteso nel bin $j$ è $\nu_j=\int_{\rm bin}[N_s^\pm S(m)+N_b^\pm B^\pm(m)]\,dm$ e il contenuto osservato $n_j$ è Poisson($\nu_j$). Il $\chi^2$ di Pearson e il pull per bin sono
\[
\chi^2=\sum_j\frac{(n_j-\nu_j)^2}{\nu_j},\qquad \text{pull}_j=\frac{n_j-\nu_j}{\sqrt{\nu_j}}.
\]
Per $k$ bin e $p$ parametri fittati ci si aspetta $\chi^2\approx k-p$ con deviazione $\sqrt{2(k-p)}$, cioè $\chi^2/{\rm ndof}=1\pm\sqrt{2/{\rm ndof}}$. Qui $k=2\times50=100$ bin da $10\MeV$, $p=12$, ndof $=88$, quindi $1\pm0.15$.

Il segnale gaussiano è un'approssimazione di una risposta del rivelatore che ha in realtà code più pesanti (scattering multiplo, radiazione, risoluzione che varia da evento a evento). Lo si controlla sostituendo la gaussiana con una \emph{Crystal Ball}: nucleo gaussiano raccordato a una coda a potenza dal lato di massa bassa.

\sez{Come si legge}
\begin{enumerate}
\item Asse $x$: $m(KKK)$; asse $y$: candidati per $10\MeV$ in \emph{scala logaritmica} (da $10$ a $\sim2000$), che rende visibile sia il fondo sia le code.
\item Elementi: punti neri con barre $\pm\sqrt n$ = dati; curva rossa = modello totale; blu tratteggiato = solo segnale; area grigia = solo fondo. Il modello è valutato alle \emph{mediane a posteriori}.
\item Pannello inferiore: pull per bin, con linee punteggiate a $\pm2$. Un buon modello ha pull che oscillano senza struttura entro $\pm2$ (con $\sim5\%$ di punti fuori).
\item Cercare \emph{raggruppamenti} di pull dello stesso segno: indicano una discrepanza sistematica, non una fluttuazione.
\end{enumerate}

\sez{Cosa dice}
\begin{enumerate}
\item \textbf{Accordo globale.} Il modello segue i dati su due ordini di grandezza, dal fondo al picco, per entrambe le cariche.
\item \textbf{Bontà numerica.} $\chi^2=95.4/50$ per $B^+$ ($1.9$ per bin) e $68.7/50$ per $B^-$ ($1.4$); totale $164.1/88=1.86$, ben oltre $1\pm0.15$. Il fit non è statisticamente perfetto.
\item \textbf{Dove sbaglia.} Un raggruppamento di pull positivi ($\approx+3$) sul fianco sinistro del picco, $5220$--$5240\MeV$, in \emph{entrambe} le cariche: i dati eccedono il modello dove la gaussiana è più bassa della coda reale. Qualche pull a $\approx5170\MeV$ è negativo.
\item \textbf{Carica per carica.} $B^+$ è peggio di $B^-$; la differenza $\chi^2(B^+)-\chi^2(B^-)=26.7$ è ripresa nella Fig.~XVI.
\item \textbf{Conseguenza per $A$.} Con una Crystal Ball $\Araw$ cambia di $-0.0002$: l'asimmetria è determinata dall'\emph{altezza relativa} dei picchi di $B^+$ e $B^-$, non dalla forma delle code; il misfit è sulle code, comune alle due cariche.
\end{enumerate}

\sez{Limiti e orale}
\begin{itemize}
\item Il $\chi^2$ binnato dipende dal binning (Fig.~XXIV): in bin da $2.5\MeV$ la stessa discrepanza dà $p=0.03$.
\item Il fit \emph{unbinned} non minimizza questo $\chi^2$, quindi la statistica non è esattamente $\chi^2_{k-p}$: la distribuzione asintotica è compresa tra $\chi^2_{k-p}$ e $\chi^2_{k-1}$ (Chernoff--Lehmann).
\item \textbf{``Se il modello non è perfetto, perché fidarsi di $\Araw$?''} Perché (i) $\Araw$ non dipende dalla forma delle code (shift $-0.0002$ con Crystal Ball), (ii) le varianti del fondo e della finestra sono nelle sistematiche, (iii) i toy di validazione dicono che stimatore e incertezze sono corretti \emph{se} il modello è giusto, e il PPC (Fig.~XV--XVI) mostra che la parte che conta per le cariche è riprodotta.
\item \textbf{``Quale modello è migliore?''} Nessuno dei provati raggiunge $\chi^2/{\rm ndof}\approx1$ (il migliore, la Crystal Ball, dà $131.0/86$): è dichiarato come limitazione.
\end{itemize}

% ======================================================================== IV
\figh{Fig.\ IV --- Posterior dei sei parametri fisici (corner plot, Sez.\ V~C)}{bayes_nominal_corner}{0.52}
\sez{Il problema}
Il risultato bayesiano è una distribuzione di probabilità sui $12$ parametri. Come la si riassume, quanto è simmetrica, e quanto $A$ dipende dagli altri parametri? Quest'ultima domanda decide se profilare e marginalizzare danno lo stesso risultato.

\sez{Teoria}
\emph{Teorema di Bayes:} $p(\theta\mid D)=L(D\mid\theta)\pi(\theta)/p(D)$. Con prior uniformi $\pi(\theta)\propto\mathbf 1[\theta\in\Theta]$ il posterior è proporzionale alla verosimiglianza dentro il box $\Theta$.

\emph{Marginalizzazione:} la distribuzione di un solo parametro si ottiene integrando sugli altri, $p(A\mid D)=\int p(\theta\mid D)\,d\theta_{\rm nuis}$; sulla catena MCMC è semplicemente l'istogramma della componente $A$. \emph{Intervallo credibile} equal-tailed al $68\%$: dal $16$° all'$84$° percentile; è una probabilità su $A$ dati i dati e il prior.

\emph{Perché il posterior è quasi gaussiano:} per $N$ grande vale il teorema di Bernstein--von Mises: $p(\theta\mid D)\to\mathcal N(\hat\theta,\,I^{-1}(\hat\theta))$, con $I_{ij}=-\partial^2\ln L/\partial\theta_i\partial\theta_j$ l'informazione di Fisher osservata, \emph{indipendentemente dal prior}. È la stessa matrice di covarianza della matrice di Hesse di \texttt{iminuit}: per questo i due metodi coincidono numericamente.

\emph{Contorni 2D:} per una gaussiana bivariata la regione a $1\sigma$ contiene $1-e^{-1/2}=39.3\%$ e quella a $2\sigma$ contiene $1-e^{-2}=86.5\%$ (non $68\%$ e $95\%$, che valgono in 1D).

\emph{Ortogonalità.} Parametrizzando $N_s^\pm=\tfrac{N_s}{2}(1\mp A)$ la \emph{somma} dei yield ($N_s$) e la loro \emph{differenza relativa} ($A$) sono quasi ortogonali: variare $A$ a $N_s$ fisso sposta candidati da una carica all'altra senza cambiarne il totale, quindi $\mathrm{Cov}(A,N_s)\approx0$.

\sez{Come si legge}
\begin{enumerate}
\item È una matrice triangolare: ogni riquadro sulla diagonale è un parametro, ogni riquadro sotto è una coppia.
\item \textbf{Diagonale:} istogramma della marginale; i tre tratti verticali sono $16$°, $50$°, $84$° percentile e i valori stanno nel titolo.
\item \textbf{Fuori diagonale:} posterior congiunto di due parametri; contorno interno $39\%$ ($1\sigma$), esterno $86\%$ ($2\sigma$).
\item \textbf{Forma = correlazione:} blob rotondo $\Rightarrow$ non correlati; ellisse inclinata $\Rightarrow$ correlati (verso dell'inclinazione = segno).
\item Confrontare la larghezza della marginale con la scala degli assi: dà l'incertezza relativa.
\end{enumerate}

\sez{Cosa dice}
\begin{enumerate}
\item \textbf{Numeri.} $N_s=15\,328^{+132}_{-133}$ (incertezza $0.9\%$), $N_b=4404^{+83}_{-81}$ ($1.9\%$), $A=-0.0443\pm0.0085$ ($19\%$ del valore, cioè $5.2\sigma$), $A_b=-0.040^{+0.018}_{-0.017}$ ($2.3\sigma$), $\mu=5284.39\pm0.16\MeV$, $\sigma=18.00\pm0.14\MeV$ ($0.8\%$).
\item \textbf{Forma.} Tutte le marginali sono simmetriche e unimodali: il posterior è quasi gaussiano, quindi mediana, moda e media coincidono entro $2\times10^{-4}$.
\item \textbf{Correlazioni.} $A$ è quasi non correlato con gli altri; la correlazione maggiore è $-0.17$ con $A_b$. L'unica ellisse visibilmente inclinata è $N_s$--$N_b$ (il totale è fissato dai dati, le due componenti si compensano).
\item \textbf{Conseguenza.} Posterior gaussiano e $A$ ortogonale ai nuisance $\Rightarrow$ marginalizzare (bayesiano) e profilare (frequentista) danno la stessa distribuzione, e Hesse, Minos e intervallo credibile coincidono (Fig.~IX).
\end{enumerate}

\sez{Limiti e orale}
\begin{itemize}
\item Sono mostrati $6$ dei $12$ parametri (gli altri sei, di forma del fondo, sono marginalizzati). La correlazione di $A$ con $A_b$ è l'unica non trascurabile: condizionare su $A_b=0$ sposta $A$ a $-0.0475$.
\item \textbf{``Perché i contorni sono al $39\%$ e $86\%$ e non al $68\%$ e $95\%$?''} Perché in due dimensioni il contenuto di $1\sigma$ è $1-e^{-1/2}$; i livelli $68\%$ e $95\%$ si ottengono con raggi $1.51\sigma$ e $2.45\sigma$ (soglie $\chi^2_2$ di $2.30$ e $5.99$, Fig.~XXVII).
\item $A_b$ è a $2.3$ deviazioni standard da zero, ma il Bayes factor lo \emph{favorisce} (Fig.~V, XIV): sono due domande diverse.
\end{itemize}

% ======================================================================== V
\figh{Fig.\ V --- Bayes factor contro l'ampiezza del prior (Sez.\ V~D)}{validation_bf_prior_width}{0.34}
\sez{Il problema}
Il $p$-value dice quanto è improbabile il dato sotto $H_0$. Il confronto bayesiano chiede una cosa diversa: quale ipotesi, $H_0$ ($A=0$) o $H_1$ ($A$ libero), il dato preferisca. La risposta bayesiana dipende dal prior: bisogna sapere \emph{quanto} dipende, perché un numero che cambia con un'arbitrarietà non è una conclusione.

\sez{Teoria}
\emph{Bayes factor:} $BF_{10}=P(D\mid H_1)/P(D\mid H_0)$, con \emph{evidenza} (verosimiglianza marginale) $P(D\mid H_i)=\int L(D\mid\theta,H_i)\,\pi(\theta\mid H_i)\,d\theta$. Moltiplicato per le probabilità a priori dà quelle a posteriori: $\frac{P(H_1\mid D)}{P(H_0\mid D)}=BF_{10}\frac{P(H_1)}{P(H_0)}$. Richiede prior \emph{propri} (integrabili).

\emph{Savage--Dickey.} Se $H_0$ è $H_1$ con $A$ fissato a $0$ e i prior sugli altri parametri sono condivisi, $\pi(\theta_{\rm nuis}\mid A=0,H_1)=\pi(\theta_{\rm nuis}\mid H_0)$, allora
\[
BF_{01}=\frac{p(A=0\mid D,H_1)}{\pi(A=0\mid H_1)}.
\]
Derivazione: $p(A=0\mid D)=\int p(D\mid A=0,\theta)\,\pi(A=0,\theta)\,d\theta/p(D)=\pi(A=0)\,P(D\mid H_0)/P(D\mid H_1)$, da cui il rapporto. Non serve integrare sui $12$ parametri: basta la marginale posteriore.

\emph{Fattore di Ockham.} Con prior uniforme su $(-W,W)$ si ha $\pi(0)=1/(2W)$, quindi $BF_{01}=2W\,p(0\mid D)$: \emph{lineare in $W$}. Un prior più largo penalizza l'ipotesi con parametro libero (spreca volume di prior): è il meccanismo del paradosso di Lindley. Numericamente, con posterior $\mathcal N(-0.0443,0.0085^2)$: $p(0\mid D)=\frac{1}{0.0085\sqrt{2\pi}}e^{-5.21^2/2}\approx6\times10^{-5}$ e $BF_{01}=2\cdot6\times10^{-5}\approx1.1\times10^{-4}$.

\emph{Scala di Jeffreys:} $BF_{10}>100$ ``decisivo'', $30$--$100$ molto forte, $10$--$30$ forte, $3$--$10$ sostanziale.

\sez{Come si legge}
\begin{enumerate}
\item Assi: $x$ = semi-ampiezza $W$ del prior uniforme $(-W,W)$ (da $0.1$ a $1$); $y$ = $BF_{01}$, entrambi in scala logaritmica.
\item Due curve: $A$ (viola continua) e $A_b$ (grigia tratteggiata). La linea punto-tratto verticale a $W=1$ è il prior nominale.
\item Orizzontali di riferimento: $BF_{01}=1$ (``evidenza uguale'': sopra, i dati favoriscono $H_0$; sotto, $H_1$) e $BF_{01}=10^{-2}$ (decisivo per $A\neq0$, $BF_{10}>100$).
\item La \emph{pendenza} $1$ di entrambe le curve in doppio-logaritmico è la linearità $BF_{01}\propto W$.
\item Cosa guardare: la curva attraversa una soglia al variare di $W$? Se sì la conclusione dipende dal prior.
\end{enumerate}

\sez{Cosa dice}
\begin{enumerate}
\item \textbf{$A$:} $BF_{01}\approx1.1\times10^{-5}$ a $W=0.1$, $1.1\times10^{-4}$ a $W=1$: una decade di $W$ costa una decade di $BF$. La curva resta sempre sotto $10^{-2}$: $H_0$ è decisivamente sfavorita per \emph{qualunque} ampiezza ragionevole. Con $W=0.3$ scende di un fattore $3.3$, conclusione immutata.
\item \textbf{$A_b$:} $BF_{01}\approx0.35$ a $W=0.1$ e $3.5$ a $W=1$. Attraversa $1$ a $W\approx0.3$: con prior stretti favorisce $A_b\neq0$ (debolmente: $BF_{10}\approx2.9$), con prior larghi favorisce $A_b=0$ (sostanziale su Jeffreys, $BF_{01}=3.5$). Nel nominale ($W=1$) i dati preferiscono $A_b=0$ \emph{nonostante} la verosimiglianza sia a $2.3\sigma$.
\item \textbf{Perché.} L'evidenza contro $A=0$ è abbastanza forte da vincere la penalità di Ockham per ogni prior; quella per $A_b\neq0$ è modesta e perde contro un prior largo.
\item \textbf{Conseguenza.} Per $A$ frequentista e bayesiano concludono lo stesso; per $A_b$ rispondono a domande diverse (un $2.3\sigma$ è ``un indizio'' al test, ma non migliora l'evidenza).
\end{enumerate}

\sez{Limiti e orale}
\begin{itemize}
\item Si usa l'approssimazione gaussiana perché a $5.2\sigma$ la catena non ha campioni a $A=0$ (Fig.~XIV).
\item \textbf{``$p=2\times10^{-7}$ e $BF_{10}=9\times10^3$ sono la stessa cosa?''} No: il $p$-value è $P(\text{dati così estremi}\mid H_0)$, non la probabilità di $H_0$, e non dipende dal prior; il Bayes factor dipende dal prior e dà il rapporto tra le verosimiglianze marginali. Qui concordano nella conclusione, non nei numeri.
\item \textbf{``Cos'è il paradosso di Lindley?''} Con un prior sempre più largo il Bayes factor favorisce $H_0$ \emph{qualunque} sia il dato: l'ipotesi con più volume di prior sprecato è penalizzata. Per questo si mostra la dipendenza da $W$.
\item \textbf{``Perché prior uniforme?''} Non informativo e con effetto trascurabile sul posterior ($N\approx2\times10^4$); conta solo per il Bayes factor.
\end{itemize}

% ======================================================================== VI
\figh{Fig.\ VI --- Potenza del test di $A=0$ (Sez.\ VII~A)}{validation_power}{0.34}
\sez{Il problema}
Un test che rigetta $A=0$ a $5.2\sigma$ è convincente solo se si conosce (i) la frequenza con cui rigetta a torto quando $A=0$ (\emph{taglia}), e (ii) la frequenza con cui rileva un effetto \emph{vero} (\emph{potenza}). Senza la prima il $p$-value non è interpretabile; senza la seconda non si sa che asimmetrie minime sarebbero risolte.

\sez{Teoria}
Nel quadro di Neyman--Pearson: \emph{taglia} $\alpha=P(\text{rigettare }H_0\mid H_0)$; \emph{potenza} $1-\beta=P(\text{rigettare }H_0\mid H_1)$, funzione di $A_{\rm true}$. La regione critica è $q(0)>c$. Per $q(0)\sim\chi^2_1$ sotto $H_0$: $c=3.84$ dà $\alpha=0.05$; $c=9$ dà $\alpha=0.0027$ ($3\sigma$ a due code).

\emph{Lemma di Neyman--Pearson:} per due ipotesi semplici il rapporto di verosimiglianza è il test più potente a taglia data; il profilo ne è l'estensione con nuisance, senza ottimalità in generale.

Nel regime gaussiano ($\hat A\sim\mathcal N(A_{\rm true},\sigma^2)$, $\sigma=0.0085$) la potenza a due code a livello $z_\alpha$ è
\[
1-\beta=\Phi\!\Big(\tfrac{|A|}{\sigma}-z_\alpha\Big)+\Phi\!\Big(-\tfrac{|A|}{\sigma}-z_\alpha\Big).
\]
Per $A=0.02$ ($|A|/\sigma=2.35$): $\alpha=0.05$ ($z=1.96$) dà $\Phi(0.39)=0.65$; $3\sigma$ ($z=3$) dà $\Phi(-0.65)=0.26$. L'asimmetria minima rivelabile con potenza $80\%$ a taglia $5\%$ è $(1.96+0.84)\sigma=2.8\sigma\approx0.024$.

\sez{Come si legge}
\begin{enumerate}
\item Asse $x$: asimmetria \emph{iniettata} nei toy, $A_{\rm true}\in\{-0.1,-0.0441,-0.02,0,+0.05\}$; asse $y$: frazione di toy in cui $q(0)$ supera la soglia, con barre binomiali.
\item Due curve: soglia $q(0)>3.84$ (taglia $5\%$, verde con cerchi) e $q(0)>9$ ($3\sigma$, grigia tratteggiata con quadrati).
\item Il punto a $A_{\rm true}=0$ è la \emph{taglia}: deve stare sulla linea punteggiata a $0.05$ (verde) e a $0.0027$ (grigia). Gli altri punti sono la \emph{potenza}: devono salire verso $1$ quando $|A_{\rm true}|$ cresce.
\item La curva è una ``V'' con minimo in $0$: potenza simmetrica attorno a zero.
\end{enumerate}

\sez{Cosa dice}
\begin{enumerate}
\item \textbf{Taglia.} $0.051\pm0.013$ per $q>3.84$ (nominale $0.05$): il test non ha falsi positivi in eccesso. Per $q>9$: $2/293=0.007$, contro $0.0027$ nominale: $2$ eventi osservati contro $0.8$ attesi, compatibile con la statistica di Poisson.
\item \textbf{Potenza.} $1.00$ a $-0.0441$, $-0.1$ e $+0.05$; $0.68$ ($0.27$ a $3\sigma$) a $-0.02$, in accordo con la previsione gaussiana ($0.65$ e $0.26$): il comportamento è quello del caso ideale.
\item \textbf{Soglia di sensibilità.} Con $19\,731$ candidati un'asimmetria di $0.02$ è rilevata circa due volte su tre; sotto $\sim0.02$ non è risolta; $80\%$ di potenza richiede $\approx0.024$. L'osservato ($-0.0441$) è ben dentro la regione di potenza $\approx1$.
\item \textbf{Conseguenza.} Il $5.2\sigma$ non è un falso positivo plausibile del test, e l'analisi non avrebbe visto un'asimmetria più piccola di $\sim0.02$: è un limite di sensibilità, non una proprietà della natura.
\end{enumerate}

\sez{Limiti e orale}
\begin{itemize}
\item Il $5.2\sigma$ dei dati ($q=27$) è un'estrapolazione della coda $\chi^2_1$: con $293$ toy a $A=0$ si verifica solo fino al livello $3\sigma$ (Fig.~XXII).
\item \textbf{``Cosa vuol dire $5.2\sigma$?''} $Z=\sqrt{q}$ con $q=27$ (convenzione a due code, $5.06\sigma$ a una coda), $p=2\times10^{-7}$; è la significatività di un'asimmetria \emph{raw}, non di violazione di CP.
\item \textbf{``Taglia e potenza si possono ottimizzare insieme?''} No: abbassare $\alpha$ riduce la potenza; il lemma di Neyman--Pearson dice quale test è il più potente a $\alpha$ fissato.
\item \textbf{``Differenza tra potenza e significatività?''} La significatività riguarda un dato osservato sotto $H_0$; la potenza riguarda la probabilità di scoprire un'asimmetria \emph{ipotetica} prima di guardare i dati.
\end{itemize}

% ======================================================================== VII
\figh{Fig.\ VII --- Mappe di Dalitz a griglia fissa $8\times8$ (Sez.\ VIII)}{bayes_nominal_dalitz_maps}{0.26}
\sez{Il problema}
L'asimmetria globale potrebbe essere la media di effetti di segno diverso concentrati in regioni cinematiche precise (la fisica delle risonanze lo suggerisce). Se così fosse, il valore globale ne nasconderebbe la struttura. Si vuole sapere se l'asimmetria è \emph{uniforme} sullo spazio delle fasi, con un test che non dipende dalla scelta della regione.

\sez{Teoria}
In un decadimento a tre corpi la cinematica ha due variabili indipendenti, p.~es.\ $m^2_{12}$ e $m^2_{13}$; la terza è fissata da $m^2_{12}+m^2_{13}+m^2_{23}=M_B^2+\sum m_i^2$. La densità dello spazio delle fasi è \emph{uniforme} nel piano $(m^2_{12},m^2_{13})$ (piano di Dalitz): $d\Gamma\propto|\mathcal M|^2\,dm^2_{12}\,dm^2_{13}$; quindi ogni struttura è dinamica (risonanze, interferenza). Le due coppie $K^+K^-$ sono equivalenti: si ordinano in $m^2_{\rm lo}\le m^2_{\rm hi}$ per eliminare l'ambiguità $H2/H3$. Una risonanza di massa $M_R$ appare come banda a $m^2=M_R^2$: $\phi(1020)$ a $1.04\GeV^2$, $D^0$ a $3.48$, $J/\psi$ a $9.59$, $\chi_{c0}$ a $11.66$.

\emph{Asimmetria locale.} In un bin con $N_i$ candidati, $A_i=(N_i^--N_i^+)/N_i$ ha incertezza binomiale $\sigma(A_i)=\sqrt{(1-A_i^2)/N_i}$, gaussiana per $N_i\gtrsim20$ (da cui il taglio a $\ge20$ candidati: $32$ dei $64$ bin). La deviazione dal valore globale è $S_i=(A_i-A_{\rm glob})/\sigma(A_i)\sim\mathcal N(0,1)$ se l'asimmetria è uniforme.

\emph{Test di omogeneità:} $\chi^2=\sum_iS_i^2$ ha $n_{\rm bin}-1=31$ gradi di libertà (uno è perso perché $A_{\rm glob}$ è stimato dai dati).

\emph{Look-elsewhere.} $P(|S|>3)=0.0027$ per bin; su $32$ bin il numero atteso è $32\times0.0027\approx0.09$, e la probabilità di due o più è $\binom{32}{2}0.0027^2\approx3\times10^{-3}$: la significatività locale va corretta per il numero di bin guardati.

\sez{Come si legge}
\begin{enumerate}
\item \textbf{(a) Statistica:} scala di grigi, candidati per bin (fino a $\approx2200$): mostra dove ci sono dati. Quasi tutto è a basso $m^2_{\rm lo}$ ($0$--$4\GeV^2$).
\item \textbf{(b) Asimmetria locale $A_i$:} colore rosso $=A_i>0$, blu $=A_i<0$, scala $\pm0.5$. Solo i bin con $\ge20$ candidati sono colorati; i bianchi sono vuoti o troppo poveri.
\item \textbf{(c) Significatività $S_i$:} stessa mappa ma normalizzata: rosso/blu intensi $=|S_i|$ grande. Riquadro nero = $|S_i|>3$.
\item \textbf{Linee tratteggiate:} $m^2$ di $D^0$, $J/\psi$, $\chi_{c0}$ in una delle due coppie (orizzontali e verticali, perché il piano è simmetrico).
\item Cosa cercare: bin colorati intensamente \emph{e} con molti candidati (in (a)): sono i più significativi.
\end{enumerate}

\sez{Cosa dice}
\begin{enumerate}
\item \textbf{Struttura.} Un gruppo di bin blu a basso $m^2_{\rm lo}$ ($0$--$4\GeV^2$) e $m^2_{\rm hi}$ tra $5$ e $12\GeV^2$; un insieme di bin rossi ad alto $m^2_{\rm hi}$ ($12$--$22\GeV^2$) con $S_i\approx+1$--$+2.5$.
\item \textbf{Bin significativi.} Due bin adiacenti a $m^2_{\rm lo}<4\GeV^2$ e $7.5<m^2_{\rm hi}<11.25\GeV^2$ hanno $S_i=-3.7$ e $-3.0$. Ne sarebbero attesi $0.09$ per fluttuazione e due o più con probabilità $3\times10^{-3}$.
\item \textbf{Test globale.} $\chi^2=80.8$ per $31$ gradi di libertà, $p=3\times10^{-6}$: l'asimmetria locale \emph{non} è uniforme. Vale anche nella sola finestra di segnale ($95.6/28$), con i veto di charm ($74.4/31$) e per griglie $6\times6$, $10\times10$, $12\times12$.
\item \textbf{Conseguenza.} L'asimmetria globale $-0.044$ è una media su una distribuzione non uniforme: negativa a basso $m^2_{\rm lo}$, più positiva ad alto $m^2_{\rm hi}$.
\end{enumerate}

\sez{Limiti e orale}
\begin{itemize}
\item Sono conteggi \emph{senza} sottrazione del fondo per pesi (sWeight) e senza correzione di rivelazione; le variazioni locali possono essere in parte strumentali, perché l'asimmetria di interazione $K^\pm$-materiale dipende dall'impulso, correlato con la posizione nel Dalitz.
\item Il $\chi^2$ globale è un test frequentista per costruzione; la controparte bayesiana è presente solo nel fit della regione pubblicata (Fig.~VIII).
\item \textbf{``Perché una griglia $8\times8$?''} Fissata \emph{prima} di guardare le asimmetrie (compromesso tra risoluzione e statistica per bin); la robustezza a $6\times6$, $10\times10$, $12\times12$ è mostrata in App.~G.
\item \textbf{``Perché non $\chi^2$ per ogni bin?''} Perché un test su $32$ bin soffre di look-elsewhere; il $\chi^2$ globale e il confronto con il valore globale lo evitano.
\end{itemize}

% ======================================================================== VIII
\figh{Fig.\ VIII --- Piano di Dalitz ordinato e regione dell'analisi pubblicata (Sez.\ VIII)}{ml_nominal_dalitz_region}{0.30}
\sez{Il problema}
Il pattern trovato (Fig.~VII) è una scoperta dei nostri dati o riproduce un risultato noto? Per rispondere senza bias di selezione si usa una regione fissata \emph{a priori} dall'articolo, non scelta guardando i nostri dati.

\sez{Teoria}
Perché l'asimmetria sia \emph{locale}: l'ampiezza di decadimento è somma di contributi, $\mathcal A=\sum_ra_re^{i(\delta_r+\phi_r)}$, con fase \emph{forte} $\delta_r$ (uguale per $B^+$ e $B^-$) e fase \emph{debole} $\phi_r$ (cambia segno nel coniugato CP). Per due ampiezze di modulo relativo $r$ e differenze di fase $\Delta\delta$, $\Delta\phi$:
\[
A_{CP}=\frac{|\bar{\mathcal A}|^2-|\mathcal A|^2}{|\bar{\mathcal A}|^2+|\mathcal A|^2}\propto \frac{-2r\sin\Delta\delta\,\sin\Delta\phi}{1+r^2+2r\cos\Delta\delta\cos\Delta\phi}.
\]
Serve una differenza di fase forte non nulla: per questo l'asimmetria si concentra dove \emph{interferiscono} risonanze, tipicamente vicino a una di esse.

L'asimmetria di \emph{conteggio} mescola segnale e fondo: $A_{\rm count}=(N_sA+N_bA_b)/(N_s+N_b)$, quindi $|A_{\rm count}|<|A|$ se il fondo è più simmetrico del segnale.

\sez{Come si legge}
\begin{enumerate}
\item \textbf{(a)} Ogni punto è un candidato nel piano $(m^2_{\rm lo},m^2_{\rm hi})$. Le linee rosse tratteggiate sono le risonanze $\phi(1020)$, $D^0$, $J/\psi$, $\chi_{c0}$ (la banda orizzontale a $m^2_{\rm hi}\approx12\GeV^2$ è la stessa $\chi_{c0}$ nell'altra coppia).
\item Il riquadro azzurro è la \emph{regione dell'analisi pubblicata}: $1.2<m^2_{\rm lo}<2.0\GeV^2$, $m^2_{\rm hi}<15\GeV^2$, appena sopra il $\phi(1020)$.
\item \textbf{(b)} Spettro di massa dei soli candidati nel riquadro: $B^+$ blu ($1273$) e $B^-$ rosso ($873$), bin da $10\MeV$.
\item Cosa guardare: la differenza di altezza tra i due picchi, e quanto fondo rimane sotto di essi.
\end{enumerate}

\sez{Cosa dice}
\begin{enumerate}
\item \textbf{Struttura.} Il piano mostra le bande delle risonanze e una zona densa a basso $m^2_{\rm lo}$. Il riquadro cade in una regione ad alta statistica.
\item \textbf{Picchi.} Nel riquadro il picco di $B^-$ ($\approx150$ per bin) è chiaramente più basso di quello di $B^+$ ($\approx250$): l'asimmetria è visibile \emph{a occhio}.
\item \textbf{Numeri.} Fit nominale ai $2146$ candidati: $\Araw=-0.234\pm0.025$ (bayesiano) e $-0.237\pm0.025$ (frequentista, $9.4\sigma$ da zero), contro $-0.239\pm0.020$ della pubblicazione. Il solo conteggio ($873$ contro $1273$) dà $\approx-0.186$: il fondo, più simmetrico, diluisce l'asimmetria; il fit, separandolo, ricava $-0.234$.
\item \textbf{Origine.} La banda del $\phi(1020)$ ($1964$ candidati) ha $-0.004\pm0.023$: l'effetto non viene dal $\phi$; i bin significativi senza di esso hanno $-0.145\pm0.019$.
\item \textbf{Conseguenza.} L'analisi riproduce, su una regione fissata a priori, un valore pubblicato indipendente: è il confronto più forte con la letteratura, e non è un risultato di CP violation (è raw).
\end{enumerate}

\sez{Limiti e orale}
\begin{itemize}
\item Campione sovrapposto, selezione e modello diversi da quelli pubblicati: la concordanza è nel valore, non una riproduzione identica.
\item \textbf{``Perché il confronto non è polarizzato?''} La regione è fissata dall'articolo, non scelta sui nostri dati; se la scegliessimo dopo aver visto la mappa avremmo un bias da selezione e bisognerebbe correggere per il look-elsewhere.
\item \textbf{``Perché serve una fase forte per un'asimmetria diretta?''} Perché $A_{CP}\propto\sin\Delta\delta\sin\Delta\phi$: senza una differenza di fase forte non c'è interferenza che distingua $B^+$ da $B^-$.
\item Non è una misura di $A_{CP}$: manca la correzione per produzione e rivelazione (App.~J).
\end{itemize}

% ======================================================================== IX
\figh{Fig.\ IX --- Profilo contro posterior; le due polarità (Sez.\ IX~A)}{validation_profile_polarity}{0.30}
\sez{Il problema}
I due metodi eliminano i nuisance in modo diverso: il frequentista \emph{profila} (massimizza), il bayesiano \emph{marginalizza} (integra). Danno la stessa distribuzione per $A$? E l'asimmetria è coerente tra le due polarità del magnete, o è un artefatto dell'una?

\sez{Teoria}
\emph{Profilo:} $L_p(A)=\max_\theta L(A,\theta)$. \emph{Marginale:} $p(A\mid D)\propto\int L(A,\theta)\pi(\theta)\,d\theta$. Se $\ln L$ è quadratica nei parametri, $L\propto\exp[-\tfrac12(x-\hat x)^TI(x-\hat x)]$, il massimo su $\theta$ e l'integrale sono entrambi proporzionali a $\exp[-(A-\hat A)^2/(2\sigma_A^2)]$ con la \emph{stessa} $\sigma_A^2=(I^{-1})_{AA}$: i due coincidono. In questo caso $\exp[-q(A)/2]=L_p(A)/L_{\max}$ normalizzata sovrappone il posterior.

\emph{Polarità.} Le tracce cariche sono curvate in direzione opposta dal campo; l'asimmetria di rivelazione destra--sinistra cambia segno invertendo la polarità. Con campioni indipendenti, la verosimiglianza congiunta è il prodotto e $q_{\rm tot}(A)=q_{\rm down}(A)+q_{\rm up}(A)$ (se i nuisance sono separati). Per stime gaussiane indipendenti la combinazione è la media pesata con l'inverso della varianza:
\[
\hat A=\frac{\hat A_1/\sigma_1^2+\hat A_2/\sigma_2^2}{1/\sigma_1^2+1/\sigma_2^2},\qquad \sigma^{-2}=\sigma_1^{-2}+\sigma_2^{-2}.
\]
Con $-0.0492\pm0.0111$ e $-0.0374\pm0.0131$ si ottiene $-0.0443\pm0.0085$, il valore del report. Compatibilità: $|\hat A_1-\hat A_2|/\sqrt{\sigma_1^2+\sigma_2^2}=0.0118/0.0172=0.69\sigma$.

\sez{Come si legge}
\begin{enumerate}
\item \textbf{(a)} Asse $x$: $\Araw$ da $-0.09$ a $0$; asse $y$: densità. L'istogramma viola è il posterior bayesiano marginale; la curva verde acqua è $\exp[-q(A)/2]$ normalizzata. Se coincidono, profilo e marginale danno la stessa distribuzione. La linea punteggiata a $0$ è l'ipotesi nulla.
\item \textbf{(b)} Asse $x$: $\Araw$; asse $y$: $q(A)=-2\Delta\ln L$ (da $0$ a $6$). Ogni curva è una parabola con minimo a $\hat A$ e larghezza $\propto\sigma$.
\item Le linee: verde = magnet down, arancione = magnet up, nera = \emph{somma} delle due $q$, tratteggiata = fit congiunto.
\item L'intervallo al $68\%$ si legge dove la parabola incontra la linea orizzontale $q=1$ (punteggiata).
\item Cosa guardare: i minimi di verde e arancione sono vicini? Nera e tratteggiata coincidono?
\end{enumerate}

\sez{Cosa dice}
\begin{enumerate}
\item \textbf{Profilo = posterior.} Nel pannello (a) le due curve sono sovrapposte: marginalizzare e profilare danno la stessa distribuzione.
\item \textbf{Polarità.} Down: $-0.0492\pm0.0111$; up: $-0.0374\pm0.0131$. I minimi sono vicini e le parabole si sovrappongono sull'intervallo centrale; la differenza è $0.7\sigma$.
\item \textbf{Somma = congiunto.} La somma $-0.0443\pm0.0084$ riproduce il fit congiunto $-0.0441\pm0.0085$: i due campioni sono indipendenti e condividono $A$, e la media pesata ricostruisce il risultato.
\item \textbf{Incertezze.} Down è più preciso di up ($0.0111$ contro $0.0131$) perché ha $1.4$ volte la statistica ($\sigma\propto1/\sqrt N$: $\sqrt{1.44}=1.2$, in accordo con $0.0131/0.0111=1.18$).
\item \textbf{Conseguenza.} Nessun indizio che l'asimmetria sia dovuta a una sola polarità; il metodo è consistente tra profilo e marginale.
\end{enumerate}

\sez{Limiti e orale}
\begin{itemize}
\item Il test copre la parte dell'asimmetria di rivelazione che cambia segno con la polarità, \emph{non} l'interazione $K^\pm$-materiale (indipendente dalla polarità).
\item \textbf{``Profilo e marginale coincidono sempre?''} No: solo se la verosimiglianza è quasi gaussiana e il prior non informativo; con forti non-linearità o prior informativi differiscono (a bassa statistica, Fig.~XXVIII, restano comunque vicini).
\item \textbf{``Perché sommare i $q$?''} Perché le due verosimiglianze sono indipendenti e il log del prodotto è la somma; se il nuisance fosse condiviso servirebbe un fit congiunto, ed è quello tratteggiato.
\item \textbf{``$0.7\sigma$ è compatibile?''} Sì: ci si aspetta di vedere differenze $>0.7\sigma$ nel $48\%$ dei casi se le polarità hanno lo stesso $A$.
\end{itemize}

% ======================================================================== X
\figh{Fig.\ X --- Cintura di confidenza di Neyman (App.\ C.2)}{validation_neyman_belt}{0.40}
\sez{Il problema}
L'intervallo del profilo ($q\le1$) è una \emph{scorciatoia} giustificata da Wilks. Qual è la costruzione \emph{definitoria} dell'intervallo di confidenza, e porta allo stesso risultato? Se sì, la scorciatoia è giustificata \emph{senza} appellarsi a un teorema asintotico.

\sez{Teoria}
\emph{Costruzione di Neyman.} Per ogni valore del parametro $\theta$ si sceglie una regione di accettazione $[x_1(\theta),x_2(\theta)]$ dello stimatore tale che $P(x\in[x_1,x_2]\mid\theta)=\mathrm{CL}$ (qui centrale, con $(1-\mathrm{CL})/2$ in ciascuna coda). L'unione delle regioni nel piano $(x,\theta)$ è la \emph{cintura} (belt). Dato $x_{\rm obs}$, l'\emph{intervallo di confidenza} è $\{\theta:x_{\rm obs}\in[x_1(\theta),x_2(\theta)]\}$, l'insieme dei valori \emph{non rifiutati}.

\emph{Proprietà:} copertura esatta per costruzione, $P(\theta_{\rm true}\in\mathrm{CI})=\mathrm{CL}$ \emph{per ogni} $\theta_{\rm true}$, perché $x_{\rm obs}$ cade nella regione di $\theta_{\rm true}$ con probabilità CL. La distribuzione di $\hat A$ dato $A_{\rm true}$ non è nota analiticamente (nuisance, Poisson): si ottiene con toy MC. L'inversione del test del rapporto di verosimiglianza, con $\hat A$ gaussiano di $\sigma$ costante, è esattamente $q(A)\le1$: è il contenuto di Wilks.

\sez{Come si legge}
\begin{enumerate}
\item Asse $y$: valore \emph{vero} $A_{\rm true}$ (da $-0.1$ a $0$, $11$ valori, $300$ toy ciascuno); asse $x$: valore \emph{stimato} $\hat A$ nei toy.
\item Le due bande: $68\%$ (scura) e $95\%$ (chiara). \textbf{Orizzontalmente}, a un $A_{\rm true}$ fissato, la banda dice dove cadrà $\hat A$ con quella probabilità.
\item La linea rossa verticale è $\hat A_{\rm obs}=-0.0441$.
\item \textbf{Verticalmente}, l'intervallo è l'insieme degli $A_{\rm true}$ la cui banda al $68\%$ contiene $\hat A_{\rm obs}$: le due linee rosse orizzontali tratteggiate sono dove la verticale esce dalla banda scura.
\item Una cintura che sale a $45^\circ$ con larghezza costante è lo stimatore ideale: non polarizzato e con $\sigma$ indipendente dal valore vero.
\end{enumerate}

\sez{Cosa dice}
\begin{enumerate}
\item \textbf{Forma.} La cintura è quasi una striscia a $45^\circ$ di larghezza costante: $\hat A$ segue $A_{\rm true}$ senza bias e con dispersione uniforme ($0.0085$) su tutto l'intervallo esplorato.
\item \textbf{Intervallo.} Il taglio dà $[-0.0528,-0.0351]$, contro $[-0.0526,-0.0356]$ del profilo (Minos): la differenza è $0.0002$ e $0.0005$, sotto $0.1\sigma$.
\item \textbf{Il $95\%$} è più largo: si legge allo stesso modo con la banda chiara.
\item \textbf{Conseguenza.} La ricetta $q\le1$ non è una convenzione ma l'intervallo di Neyman: stessa copertura per costruzione, ottenuta qui con $3300$ toy.
\end{enumerate}

\sez{Limiti e orale}
\begin{itemize}
\item Costruita per $A_{\rm true}\in[-0.1,0]$ con interpolazione lineare tra $11$ punti: i bordi sono a $\hat A\approx-0.12$ e $+0.02$.
\item Feldman--Cousins (ordinamento per rapporto di verosimiglianza) serve vicino a un bordo fisico del parametro; qui no: $A\in(-1,1)$ e il valore stimato è $5\sigma$ da $0$.
\item \textbf{``Cos'è un intervallo di confidenza al $68\%$?''} Non la probabilità che $A$ sia nell'intervallo osservato, ma la frazione di esperimenti ripetuti in cui l'intervallo costruito contiene il valore vero.
\item \textbf{``Un intervallo di confidenza può essere vuoto o contenere valori impossibili?''} Sì, in casi limite; qui no perché $\hat A$ è lontano dai bordi.
\end{itemize}

% ======================================================================== XI

\figh{Fig.\ XI --- $\Araw$ per terzili del momento del $B$ (App.\ C.2)}{momentum_terciles}{0.30}

\sez{Il problema}
L'asimmetria di rivelazione dei kaoni dipende dal loro impulso (le sezioni d'urto di $K^+$ e $K^-$ nel materiale cambiano con $p$), e l'asimmetria di produzione del $B$ dipende da $p_T$ e pseudorapidità. Se $\Araw$ fosse in parte strumentale o di produzione dovrebbe variare con il momento del $B$. Si misura in tre intervalli.

\sez{Teoria}
La divisione in \emph{terzili} produce tre sottocampioni indipendenti con lo stesso numero di candidati ($6577$ ciascuno), così che le tre misure abbiano precisione confrontabile e nessun bin sia statisticamente povero. Per tre stime indipendenti $\hat A_k\pm\sigma_k$ la compatibilità con un valore comune si valuta con
\[
\chi^2=\sum_k\frac{(\hat A_k-\bar A)^2}{\sigma_k^2},\qquad \bar A=\frac{\sum_k\hat A_k/\sigma_k^2}{\sum_k1/\sigma_k^2},
\]
distribuito come $\chi^2$ con $k-1=2$ gradi di libertà sotto l'ipotesi di valore comune. Un trend con il momento darebbe $\chi^2$ grande e un andamento monotòno.

\sez{Come si legge}
\begin{enumerate}
\item Pannello (a): terzili di $p_T$ del $B$ (in GeV: $1$--$5$, $5$--$8$, $8$--$31$); pannello (b): terzili di $p$ ($27$--$78$, $78$--$110$, $110$--$289$). Asse $y$: $\Araw$ del fit frequentista congiunto.
\item Punti con barre = stima per terzile; fascia grigia = valore globale $\pm1\sigma$; linea punteggiata = valore globale.
\item Cosa guardare: i punti stanno nella fascia? Hanno un andamento monotòno con il momento?
\end{enumerate}

\sez{Cosa dice}
\begin{enumerate}
\item \textbf{$p_T$:} $-0.0563\pm0.0169$, $-0.0362\pm0.0141$, $-0.0452\pm0.0134$. Il primo e il secondo differiscono di $1.4\sigma$ (la massima deviazione citata nel testo); non c'è monotonia: il valore centrale non è né il più basso né il più alto.
\item \textbf{$p$:} $-0.0391\pm0.0139$, $-0.0475\pm0.0149$, $-0.0478\pm0.0153$: quasi piatto.
\item \textbf{Compatibilità.} Valore comune $-0.0447$ ($p_T$) e $-0.0445$ ($p$), $\chi^2=0.83$ e $0.24$ per $2$ gradi di libertà ($p=0.66$ e $0.89$): perfettamente compatibili.
\item \textbf{Conseguenza.} Nessuna dipendenza evidente dal momento: $\Araw$ non mostra la firma di un effetto strumentale legato al momento, e un'eventuale correzione di $A_D(p)$ non cambierebbe il quadro al livello della precisione.
\end{enumerate}

\sez{Limiti e orale}
\begin{itemize}
\item Tre intervalli di $\sim6600$ candidati hanno ciascuno $\sigma\approx0.014$--$0.017$: sono sensibili solo a variazioni $\gtrsim0.03$. La mancanza di trend non esclude effetti più piccoli.
\item \textbf{``Perché i terzili e non intervalli fissi?''} Perché a numero uguale di candidati le tre misure hanno la stessa precisione.
\item \textbf{``È una correzione?''} No: è un controllo. Una correzione di $A_D(p)$ richiede la misura nel canale di controllo (App.~J).
\end{itemize}


\figh{Fig.\ XII --- Incertezze dei dodici parametri con tre metodi (App.\ C.3)}{uncertainties_12}{0.30}

\sez{Il problema}
L'incertezza di Hesse è un'approssimazione asintotica (parabola di $\ln L$ al minimo). Dà la stessa incertezza di due stime indipendenti, bootstrap sui dati e posterior, non solo per $A$ ma per \emph{tutti} i parametri, compresi i nuisance?

\sez{Teoria}
\emph{Hesse:} $\mathrm{Cov}(\hat\theta)\simeq I^{-1}$ con $I_{ij}=-\partial^2\ln L/\partial\theta_i\partial\theta_j$ al massimo. \emph{Bootstrap:} la distribuzione campionaria dello stimatore si approssima rifittando $B=300$ ricampionamenti con reinserimento dei dati; la deviazione standard delle stime è l'incertezza, senza assumere una parabola. \emph{Posterior:} deviazione standard della marginale nella catena.

\emph{Caso dei yield.} Il ricampionamento dei $19\,731$ candidati ne mantiene fisso il \emph{totale}: $N_s+N_b=N$ costante, quindi $N_s$ e $N_b$ sono perfettamente anticorrelati con $\sigma(N_s)=\sigma(N_b)$ (qui $72.5$ entrambi). Nel fit esteso invece il totale è Poisson ($\sqrt{N}\approx140$), che l'Hesse include ($\sigma(N_s)=133$, $\sigma(N_b)=82$). La differenza è dunque attesa, non un difetto.

\sez{Come si legge}
\begin{enumerate}
\item Asse $x$: i dodici parametri del fit ($A$, $A_b$, $N_s$, $N_b$, $\mu$, $\sigma$ e i parametri di forma del fondo per $B^\pm$). Asse $y$: rapporto $\sigma_{\rm metodo}/\sigma_{\rm Hesse}$.
\item Quadrati verdi = bootstrap dei dati; cerchi viola = posterior. Linea nera a $1$ (accordo perfetto), fascia grigia $\pm10\%$.
\item Cosa guardare: i punti stanno nella fascia? Quali escono, e perché?
\end{enumerate}

\sez{Cosa dice}
\begin{enumerate}
\item \textbf{Posterior.} Dentro il $10\%$ per tutti i parametri (rapporti $0.96$--$1.09$): coincide con Hesse, come atteso per un posterior quasi gaussiano.
\item \textbf{Bootstrap per i parametri fisici.} $A$ ($0.997$), $A_b$ ($1.035$), $\mu$ ($1.00$), $\sigma$ ($1.07$): nella fascia.
\item \textbf{Eccezioni.} $N_s$ ($0.55$) e $N_b$ ($0.88$): è l'effetto del totale fisso. Anche $\lambda_1^+$ ($0.74$): la pendenza ripida del fondo $B^+$ è il parametro peggio vincolato (cfr.~il trace plot, Fig.~XIII) e la sua verosimiglianza è non parabolica.
\item \textbf{Conseguenza.} Per $A$ e per i parametri che contano le tre stime coincidono: l'incertezza di Hesse è affidabile.
\end{enumerate}

\sez{Limiti e orale}
\begin{itemize}
\item \textbf{Da correggere nel testo.} Il testo dell'App.~C.3 dice che i tre metodi concordano ``entro il 10\% per tutti i parametri di forma'': per $\lambda_1^+$ il bootstrap differisce del $26\%$. La frase va attenuata (``tranne $N_s$, $N_b$ e $\lambda_1^+$'').
\item \textbf{``Quale usi per $A$?''} Hesse è quella riportata; Minos (non parabolica) e bootstrap la confermano.
\item \textbf{``Perché il bootstrap dà $\sigma(N_s)=\sigma(N_b)$?''} Perché il totale è fisso, quindi si compensano esattamente.
\end{itemize}

\figh{Fig.\ XIII --- Tracce dei walker (App.\ D.5)}{bayes_nominal_trace}{0.52}
\sez{Il problema}
Una catena MCMC campiona il posterior solo dopo essere \emph{convergita}. Se i walker partono da punti diversi ed arrivano alla stessa regione, e restano stazionari, la catena è affidabile. Bisogna \emph{vedere} che non è intrappolata vicino al punto di partenza né in una moda secondaria.

\sez{Teoria}
\emph{MCMC:} si costruisce una catena di Markov con distribuzione stazionaria $p(\theta\mid D)$; dopo un transitorio (\emph{burn-in}) i campioni sono estrazioni (correlate) dal posterior. \emph{emcee} usa il campionatore d'insieme affine-invariante di Goodman e Weare: $K$ walker; per ognuno si propone $Y=X_j+Z\,(X_k-X_j)$ con $X_k$ un altro walker, $Z$ con densità $g(z)\propto1/\sqrt z$ su $[1/a,a]$ ($a=2$), e si accetta con probabilità $\min\{1,Z^{d-1}p(Y)/p(X_k)\}$, $d$ = dimensione. L'invarianza affine lo rende insensibile alle correlazioni e alle diverse scale dei parametri.

\emph{Autocorrelazione.} Campioni consecutivi sono correlati: $\rho(t)=\mathrm{Corr}(\theta_s,\theta_{s+t})$. Il tempo di autocorrelazione integrato è $\tau=1+2\sum_{t\ge1}\rho(t)$ e il numero di campioni indipendenti è $N_{\rm eff}\approx N_{\rm samples}/\tau$. Una stima di $\tau$ è affidabile se la catena è lunga $\gtrsim50\tau$.

\sez{Come si legge}
\begin{enumerate}
\item Quattro righe: $A$, $N_s$, $\sigma$, frazione di fondo ripido $f_+$. Ogni linea colorata è un walker ($32$), registrato ogni $10$ passi, partiti uniformemente sull'\emph{intera} scatola dei prior (non vicino al massimo).
\item \textbf{Colonna sinistra:} primi $3000$ passi, scala completa del prior (per $A$: $-1$ a $1$). Si vede il transitorio.
\item \textbf{Colonna destra:} i $25\,000$ passi, ingranditi sul posterior. La linea punteggiata verticale è il burn-in (qui la prima metà).
\item Un buon risultato: il transitorio finisce presto, poi \emph{tutte} le linee si sovrappongono in una striscia stazionaria (un ``bruco'') senza deriva, senza walker isolati e senza salti tra livelli.
\end{enumerate}

\sez{Cosa dice}
\begin{enumerate}
\item \textbf{Convergenza.} Tutti i walker raggiungono il bulk entro $\sim2000$ passi: $A$ da $[-1,1]$ a $\approx-0.044$, $N_s$ da $\sim75\,000$ a $\approx15\,300$, $\sigma$ da $\sim35$ a $18\MeV$.
\item \textbf{Il più lento è $f_+$.} Mostra un'escursione ($\approx0.5$ attorno al passo $1800$) prima di assestarsi a $\approx0.1$, ben dentro la metà scartata: la forma del fondo è il parametro peggio vincolato.
\item \textbf{Stazionarietà.} Dopo il burn-in le strisce sono costanti per $\sim23\,000$ passi, senza trend.
\item \textbf{Unica moda.} Due run di questo tipo danno $\Araw=-0.0446$ e $-0.0442$, come il nominale.
\item \textbf{Diagnostica nominale.} $\tau\in[147,183]$ passi, $N/\tau\ge113$, $\gtrsim3600$ campioni effettivi per parametro.
\end{enumerate}

\sez{Limiti e orale}
\begin{itemize}
\item Una catena non può \emph{provare} la convergenza, solo non rilevare problemi: per questo si usano più run da punti di partenza dispersi (e il tempo di autocorrelazione).
\item \textbf{``Cosa sono label switching e multimodalità?''} Le due esponenziali del fondo sono scambiabili (se non si imponesse un ordine), il che crea due modi equivalenti; l'ordinamento $\lambda_2<\lambda_1$ nel prior elimina la simmetria.
\item \textbf{``Perché $32$ walker?''} Per emcee servono almeno $2d=24$ walker per $d=12$ parametri; $32$ è un margine.
\item \textbf{``Come scegli il burn-in?''} Quando i walker sono tutti dentro la regione stazionaria nel grafico e $\tau$ è stato stimato; qui $1/6$ della catena nominale (quella con avvio vicino al massimo) e metà in questo test, più prudente.
\end{itemize}

% ======================================================================== XII
\figh{Fig.\ XIV --- Rapporto di Savage--Dickey (App.\ D.11)}{bayes_nominal_savage_dickey}{0.30}
\sez{Il problema}
La Fig.~V mostra il Bayes factor in funzione del prior; qui si vede \emph{come} è costruito dai suoi due ingredienti: la densità posteriore e la densità del prior nel punto nullo. Permette anche di giudicare se l'approssimazione gaussiana in coda è ragionevole.

\sez{Teoria}
$BF_{01}=p(\theta=0\mid D)/\pi(\theta=0)$ (derivazione in Fig.~V). Il prior su $A$, $A_b$ è uniforme su $(-1,1)$ con densità $0.5$. La densità posteriore a $0$ va stimata in una \emph{coda}: a $5.2\sigma$ la probabilità oltre $0$ è $\sim2\times10^{-7}$, quindi in $6.7\times10^5$ campioni se ne attendono $\sim0.13$: la catena non contiene punti a $A=0$ e un istogramma o un kernel estimator non servono (rifletterebbero la banda del kernel). Si usa l'approssimazione gaussiana con media e deviazione standard della catena, giustificata dalla forma del posterior (Fig.~IV).

\sez{Come si legge}
\begin{enumerate}
\item Scala \emph{logaritmica} sull'asse $y$ (densità, da $10^{-5}$ a $10^2$); asse $x$: il parametro ($A$ in (a), $A_b$ in (b)).
\item Istogramma viola = marginale della catena; curva rossa = approssimazione gaussiana; linea tratteggiata orizzontale = prior ($0.5$).
\item Sulla verticale in $0$ (punteggiata) due punti: grigio = densità del prior ($0.5$), rosso = densità posteriore (gaussiana) a $0$.
\item $BF_{01}$ è il rapporto \emph{posteriore/prior} a $0$, quindi la distanza verticale (in dex) tra i due punti. Se il punto rosso sta \emph{sotto} il grigio, $H_0$ è sfavorita.
\end{enumerate}

\sez{Cosa dice}
\begin{enumerate}
\item \textbf{(a) $A$:} densità posteriore a $0$ $\approx5.5\times10^{-5}$ contro $0.5$ del prior: $BF_{01}\approx1.1\times10^{-4}$, i dati abbassano la densità di $\sim4$ ordini di grandezza.
\item \textbf{Coda vuota.} L'istogramma della catena ha bin vuoti oltre $A\gtrsim-0.012$: la coda non è campionata, e per questo si usa la gaussiana.
\item \textbf{(b) $A_b$:} a $2.3\sigma$ da zero ha densità posteriore a $0$ \emph{sopra} il prior ($\approx1.7$ contro $0.5$): $BF_{01}\approx3.5$.
\item \textbf{Perché sopra il prior.} La verosimiglianza ha un massimo intorno a $-0.04$ ma è larga ($\pm0.018$): il posterior si concentra in $\sim0.07$ di larghezza, quindi la densità è $\sim1/0.07\approx14$ al picco e ancora $\approx1.7$ a $0$, mentre il prior è sparso su $2$ unità ($0.5$). Il guadagno di verosimiglianza non compensa lo spargimento del prior: Ockham.
\item \textbf{Conseguenza.} I dati favoriscono mildly $A_b=0$ pur essendo $A_b$ a $2.3\sigma$: significatività e evidenza sono cose diverse.
\end{enumerate}

\sez{Limiti e orale}
\begin{itemize}
\item La scala logaritmica ingigantisce le code: la gaussiana è accurata dove c'è massa, la sua coda a $5\sigma$ è un'estrapolazione (controllata solo dal fatto che la forma è gaussiana al centro).
\item \textbf{``Perché non usare un kernel density estimator?''} Perché senza campioni nella coda restituisce solo la banda del kernel.
\item \textbf{``Savage--Dickey vale sempre?''} Solo per ipotesi annidate con prior condizionali condivisi; altrimenti serve l'evidenza (bridge sampling o simili).
\item \textbf{``Sensibilità al prior?''} $BF\propto W$; vedi Fig.~V.
\end{itemize}

% ======================================================================== XIII
\figh{Fig.\ XV --- Posterior predictive check (App.\ D.13)}{bayes_nominal_ppc}{0.22}
\sez{Il problema}
Il $\chi^2$ calcolato alle mediane a posteriori usa un solo valore dei parametri e ignora l'incertezza su di essi. Come si valuta l'adeguatezza del modello \emph{tenendone conto}, e con statistiche mirate a caratteristiche specifiche dei dati?

\sez{Teoria}
\emph{Distribuzione predittiva a posteriori:}
\[
p(y^{\rm rep}\mid y)=\int p(y^{\rm rep}\mid\theta)\,p(\theta\mid y)\,d\theta.
\]
In pratica: per $j=1,\dots,1000$ si estrae $\theta_j$ dalla catena, si genera $y^{\rm rep}_j$ dal modello, e si calcola una \emph{discrepanza} $T(y,\theta)$ sia sulla replica sia sui dati. Il $p$-value bayesiano è $p_B=P\big(T(y^{\rm rep},\theta)\ge T(y,\theta)\mid y\big)$. Non è uniforme in generale (usa i dati due volte, per il fit e per il test): valori molto vicini a $0$ o $1$ segnalano un problema; per statistiche \emph{fittate} ci si attende $p_B$ vicino a $0.5$.

\sez{Come si legge}
\begin{enumerate}
\item Tre pannelli, ognuno con una discrepanza diversa sull'asse $x$; l'asse $y$ è il numero di repliche.
\item L'istogramma azzurro è la distribuzione della discrepanza nelle $1000$ repliche; la linea rossa verticale è il valore osservato nei dati (mediana su $\theta_j$).
\item $p_B$ è la frazione di istogramma \emph{a destra} della linea rossa (nei pannelli (a); in (b) e (c) a due code). La legenda riporta il valore.
\item Cosa guardare: la linea rossa cade dentro l'istogramma (modello adeguato per quella statistica) o fuori (inadeguato)?
\end{enumerate}

\sez{Cosa dice}
\begin{enumerate}
\item \textbf{(a) $\chi^2$ binnato.} Repliche centrate a $\approx104$ (larghezza $15$), osservato $175$ ben a destra: $p_B<10^{-3}$. Nessuna replica raggiunge i dati: il modello \emph{non riproduce la forma} dello spettro. Il misfit è concentrato in $B^+$ ($p_B<10^{-3}$ contro $0.03$ per $B^-$).
\item \textbf{(b) Asimmetria della finestra.} $p_B=0.996$: il dato è riprodotto, ma è \emph{ovvio} (vedi limiti).
\item \textbf{(c) Asimmetria delle bande.} Replice centrate a $\approx-0.01$, osservato $\approx-0.01$, $p_B=0.73$: l'asimmetria nelle bande, che il fit vincola meno direttamente, è riprodotta.
\item \textbf{Conseguenza.} Il modello descrive male la \emph{forma} (code del segnale) ma non mostra distorsioni dipendenti dalla carica nelle bande.
\end{enumerate}

\sez{Limiti e orale}
\begin{itemize}
\item (b) \emph{non è informativo}: $A$ è un parametro fittato, quindi le repliche sono centrate sul dato per costruzione. Per questo la Fig.~XVI usa statistiche non fittate.
\item \textbf{``Differenza tra $p_B$ e il $p$-value classico?''} $p_B$ marginalizza sul posterior dei parametri (i parametri sono incerti), il $p$-value classico fissa $\theta$ (ipotesi nulla). $p_B$ non è uniforme sotto il modello corretto; è conservativo.
\item \textbf{``Perché il PPC e non solo il $\chi^2$?''} Perché testa specifiche caratteristiche (asimmetria nelle bande) con l'incertezza dei parametri.
\item Il $\chi^2$ in (a) è calcolato \emph{con i parametri della replica}, non alle mediane: per questo $\approx104$ e non $\approx88$.
\end{itemize}

% ======================================================================== XIV
\figh{Fig.\ XVI --- PPC con statistiche non fittate dal modello (App.\ D.14)}{validation_ppc_unfitted}{0.22}
\sez{Il problema}
Un controllo ha potere solo se testa qualcosa che il fit \emph{non} ha già aggiustato. La Fig.~XV (b) è banale perché $A$ è fittato. Quali statistiche sono \emph{predette} dal modello senza essere adattate, e che dicono delle cariche?

\sez{Teoria}
In una fetta di massa $m$ l'asimmetria di conteggio è la media pesata con la frazione locale di segnale:
\[
A_s(m)=\frac{N_s(m)\,A+\big[N_b^-(m)-N_b^+(m)\big]}{N_s(m)+N_b(m)}\ \to\ \text{al picco }\simeq A,\ \text{nei lati }\simeq A_b(m).
\]
Poiché il fondo di ciascuna carica ha una forma propria ($\lambda^\pm$, $f_\pm$), $A_b(m)$ \emph{dipende dalla massa}: il modello predice un andamento a gradini non banale che non è stato fittato. Statistica sintetica: $T=\sum_s(A_s-\bar A_s)^2/\sigma_s^2$ (circa $\chi^2$ con $\sim10$ gradi di libertà). Seconda statistica: $\chi^2(B^+)-\chi^2(B^-)$, che confronta la qualità del fit tra le cariche.

\sez{Come si legge}
\begin{enumerate}
\item \textbf{(a)} Asse $x$: massa; asse $y$: asimmetria di conteggio in fette da $50\MeV$. Punti neri con barre = dati; \emph{gradino rosso} = previsione del modello (mediane a posteriori). Se il modello è giusto i punti seguono il gradino.
\item \textbf{(b)} Distribuzione di $T$ nelle repliche (istogramma) e valore osservato (linea rossa); $p_B$ in legenda.
\item \textbf{(c)} Distribuzione di $\chi^2(B^+)-\chi^2(B^-)$ nelle repliche, centrata su $0$ per costruzione, e valore osservato.
\item Cosa guardare: i punti in (a) si discostano dal gradino in modo sistematico? La linea rossa in (b), (c) è lontana dal centro?
\end{enumerate}

\sez{Cosa dice}
\begin{enumerate}
\item \textbf{(a) Previsione corretta.} Il modello prevede $\approx-0.045$ sotto il picco, circa $-0.07$ nelle fette a bassa massa e valori positivi ($\approx+0.05$) ad alta massa; i dati seguono questo andamento a gradini entro le barre (a parte uno scarto isolato a $\approx5175\MeV$).
\item \textbf{(b)} Valore osservato $\approx7.5$, repliche centrate a $\approx10$: $p_B=0.67$. La struttura in massa dell'asimmetria è riprodotta.
\item \textbf{(c)} $27$ nei dati contro $-0\pm15$ nelle repliche: $p_B=0.08$, $1.8\sigma$. Il campione $B^+$ è peggio descritto di $B^-$.
\item \textbf{Conseguenza.} È l'unico controllo bayesiano con potere sulle cariche: nessuna distorsione carica-dipendente evidente, e il misfit sta in $B^+$ in modo mild.
\end{enumerate}

\sez{Limiti e orale}
\begin{itemize}
\item Un buon PPC \emph{non} prova che il modello sia giusto, solo che non sbaglia in quelle direzioni.
\item \textbf{``Perché l'asimmetria nelle fette cambia con la massa?''} Perché segnale e fondo hanno asimmetrie diverse e le fette sono dominate da uno o dall'altro; l'asimmetria del fondo dipende dalla massa perché i due fondi hanno pendenze diverse.
\item \textbf{``È un test o un grafico?''} Ha una statistica ($T$) e un $p_B$: è un test, e il grafico (a) ne è l'illustrazione.
\item Con $\sim10$ fette il test ha poca potenza su deviazioni piccole.
\end{itemize}

% ======================================================================== XV
\figh{Fig.\ XVII --- Posterior di $\Araw$ per polarità (App.\ D.15)}{bayes_nominal_magnet}{0.33}
\sez{Il problema}
L'asimmetria potrebbe essere un artefatto di una sola polarità del magnete (effetto strumentale destra--sinistra). Si misura separatamente e si guarda se le due misure concordano.

\sez{Teoria}
\emph{Asimmetria di rivelazione.} $A_D$ nasce da due effetti: (i)~$K^+$ e $K^-$ hanno sezioni d'urto nucleari diverse (il $K^-$ interagisce di più col materiale), un effetto che dipende dall'impulso e \emph{non} dalla polarità; (ii)~la curvatura nel campo è opposta per le due cariche, e tracce che escono dall'accettanza a destra o a sinistra del rivelatore sono perse in modo diverso: un effetto \emph{destra--sinistra} che \emph{cambia segno} invertendo la polarità. Per $A_{\rm CP}$ fisico, invece, il segno resta lo stesso. Quindi un'asimmetria prodotta dall'effetto (ii) apparirebbe con segni \emph{opposti} nei due campioni e si cancellerebbe nella media. Le due misure sono indipendenti; sono compatibili se $|\hat A_{\rm down}-\hat A_{\rm up}|\lesssim\sqrt{\sigma^2_{\rm down}+\sigma^2_{\rm up}}$.

\sez{Come si legge}
\begin{enumerate}
\item Asse $x$: $\Araw$; asse $y$: densità posteriore normalizzata. Tre istogrammi: magnet down (verde), magnet up (arancione), campione completo (nero).
\item La linea grigia punteggiata verticale a $0$ è l'ipotesi nulla.
\item La \emph{larghezza} di ciascun istogramma è l'incertezza; la \emph{posizione} del picco è la stima.
\item Cosa guardare: i due istogrammi colorati si sovrappongono molto? Il nero cade tra i due (media pesata) ed è più stretto? Sono tutti a sinistra dello zero?
\end{enumerate}

\sez{Cosa dice}
\begin{enumerate}
\item \textbf{Down:} $-0.0497\pm0.0112$ ($n=11\,636$). \textbf{Up:} $-0.0378\pm0.0129$ ($n=8095$). \textbf{Combinato:} $-0.0443\pm0.0085$ ($n=19\,731$).
\item \textbf{Concordanza.} I picchi stanno a $-0.050$ e $-0.038$; le code si sovrappongono sull'intervallo $[-0.06,-0.03]$. Differenza $0.0119$, cioè $0.7\sigma$.
\item \textbf{Segno.} Entrambi sono negativi e lontani da $0$ ($4.4\sigma$ e $2.9\sigma$ dallo zero): nessuna polarità ha $A\approx0$.
\item \textbf{Combinato più stretto.} La larghezza $0.0085$ è $\approx\sqrt{0.0112^{-2}+0.0129^{-2}}^{-1}$: l'incertezza scala come previsto.
\item \textbf{Conseguenza.} Un artefatto prodotto dall'effetto (ii) avrebbe dato segni opposti; non c'è.
\end{enumerate}

\sez{Limiti e orale}
\begin{itemize}
\item Non testa l'effetto (i), indipendente dalla polarità: per questo $A_{\rm raw}\neq A_{CP}$ e servirebbe un canale di controllo (App.~J).
\item \textbf{``Perché le due polarità non danno lo stesso valore?''} Non devono: sono campioni indipendenti con fluttuazioni statistiche; $0.7\sigma$ di differenza è nella norma.
\item \textbf{``Perché non si fissa $A_{\rm down}=A_{\rm up}$?''} Lo si fa nel fit congiunto (nominale); qui si verifica che l'assunzione sia ragionevole.
\end{itemize}

% ======================================================================== XVI
\figh{Fig.\ XVIII --- Copertura del posterior su toy MC (App.\ D.16)}{bayes_nominal_toys}{0.30}
\sez{Il problema}
L'intervallo credibile è una probabilità sul parametro, \emph{non} ha copertura garantita. Con prior piatti e $N\approx2\times10^4$ la ha in pratica? Se sì, è legittimo confrontarlo con quello frequentista.

\sez{Teoria}
\emph{Bootstrap parametrico:} si genera da $\theta=\hat\theta$ (mediane a posteriori) con yield di Poisson per carica e componente; le masse gaussiane ed esponenziali troncate si estraggono per inversione della funzione di ripartizione (metodo della trasformazione, $m=F^{-1}(u)$, $u\sim U(0,1)$); si rifà il fit con una catena di $12\,000$ passi. Per ciascuna replica si registra se l'intervallo contiene $A_{\rm true}$. La frazione ha incertezza binomiale $\sqrt{p(1-p)/n}=\sqrt{0.68\cdot0.32/200}=0.033$. Il quadro è quello dell'inferenza bayesiana \emph{calibrata} (valutare con criteri frequentisti una procedura bayesiana): garantita asintoticamente dal teorema di Bernstein--von Mises, da verificare altrimenti.

\sez{Come si legge}
\begin{enumerate}
\item \textbf{(a)} Pull per replica: $(\hat A-A_{\rm true})/\sigma(\hat A)$ con $\hat A$ mediana e $\sigma$ deviazione standard del posterior. Istogramma azzurro contro $\mathcal N(0,1)$ (rossa); in alto media e deviazione standard.
\item \textbf{(b)} Per i livelli $68\%$ e $95\%$: punto nero con barre binomiali = frazione osservata di repliche il cui intervallo contiene $A_{\rm true}$; tratto rosso = valore nominale.
\item Cosa guardare: il punto nero è compatibile col tratto rosso entro le barre? Il pull è centrato su $0$ con larghezza $1$?
\end{enumerate}

\sez{Cosa dice}
\begin{enumerate}
\item \textbf{Stimatore.} Pull con media $-0.03\pm0.07$ e deviazione standard $1.04$; la media delle stime differisce dal valore iniettato di $-0.0002\pm0.0006$ e la loro dispersione, $0.0088$, concorda con la deviazione standard del posterior ($0.0085$).
\item \textbf{Copertura al $68\%$:} $0.650\pm0.033$, a $1.0\sigma$ sotto il nominale $0.683$. \textbf{Al $95\%$:} $0.930\pm0.015$, a $1.3\sigma$ sotto $0.954$.
\item \textbf{Direzione.} Entrambe leggermente \emph{sotto} il nominale, ma entro $1.3\sigma$: compatibili. Con $200$ toy una sotto-copertura di $3$ punti percentuali non è distinguibile dalla fluttuazione.
\item \textbf{Conseguenza.} Con prior piatti e questa statistica l'intervallo credibile ha, in pratica, la copertura nominale: sul piano numerico i due metodi sono intercambiabili.
\end{enumerate}

\sez{Limiti e orale}
\begin{itemize}
\item È un test di \emph{chiusura} (dati generati dal modello), diverso dal bootstrap sui dati della Fig.~II: testa la procedura, non il modello.
\item La legenda dice ``pseudo-datasets'': il report usa ``toy MC''.
\item $200$ toy bastano solo per controlli a $\sim3\%$; con $1000$ (come il frequentista) l'errore scenderebbe a $1.5\%$.
\item \textbf{``Perché una catena di $12\,000$ passi per toy e non $25\,000$?''} Per limitare il tempo di calcolo ($200$ fit); $\tau\approx150$ passi, quindi $12\,000$ sono $\sim80\tau$, sopra la soglia di $50$.
\end{itemize}


% ======================================================================== NEW syst_variants
\figh{Fig.\ XIX --- Spostamenti di tutte le varianti (App.\ E.5)}{syst_variants}{0.50}
\sez{Il problema}
La Tab.~II dà tredici spostamenti per due metodi: ventisei numeri, più le incertezze attese. Servono due cose che la tabella mostra male: (i) i due metodi danno gli stessi spostamenti \emph{variante per variante}; (ii) quale gruppo domina, e quanto di ogni spostamento è compatibile con una fluttuazione statistica.

\sez{Teoria}
Lo studio delle sistematiche cambia una scelta alla volta rispetto al nominale e misura lo spostamento $\Delta=A^{\rm var}_{\rm raw}-A^{\rm nom}_{\rm raw}$. Le regole sono quattro.
\begin{enumerate}
\item \textbf{Varianti che cambiano il campione} (finestra, preselezione): il campione variato è contenuto nel nominale, quindi parte di $\Delta$ è una fluttuazione statistica di dimensione attesa $\sigma_B=\sqrt{|\sigma_{\rm var}^2-\sigma_{\rm nom}^2|}$ (Fig.~XXI). La variante contribuisce solo se $|\Delta|>\sigma_B$.
\item \textbf{Varianti sullo stesso campione} (fondo, segnale): non c'è fluttuazione statistica da sottrarre, si usa $\Delta$ intero.
\item \textbf{Varianti rifiutate dai dati} ($\Delta\chi^2>25$ sugli stessi dati): escluse; è il caso dell'esponenziale singola ($\Delta\chi^2=134$).
\item \textbf{Combinazione:} per gruppo (finestra, fondo, segnale, preselezione) si prende il massimo contributo; i gruppi si sommano in quadratura.
\end{enumerate}
Il raggruppamento riflette che varianti dello stesso gruppo sono \emph{alternative} per la stessa scelta (non effetti indipendenti), mentre gruppi diversi sono scelte indipendenti.

\sez{Come si legge}
\begin{enumerate}
\item Asse $y$: le $13$ varianti, in quattro gruppi separati da linee punteggiate (finestra, fondo, segnale, preselezione). Asse $x$: $\Delta$.
\item Per ogni variante due punti: bayesiano (viola, cerchio) e frequentista (verde, quadrato), leggermente sfalsati verticalmente.
\item Le barre orizzontali sono $\pm\sigma_B$, presenti solo per le varianti che cambiano il campione; le varianti sullo stesso campione non hanno barra.
\item Simbolo vuoto = variante rifiutata dai dati. La linea verticale è $\Delta=0$.
\item Cosa guardare: (i) cerchio e quadrato coincidono? (ii) la barra $\sigma_B$ contiene lo zero (spostamento compatibile con fluttuazione)? (iii) quale gruppo ha gli spostamenti più grandi?
\end{enumerate}

\sez{Cosa dice}
\begin{enumerate}
\item \textbf{I due metodi coincidono.} Per ogni variante le posizioni bayesiana e frequentista differiscono per meno di $0.0005$: stesse scelte, stessa verosimiglianza, stessa regola.
\item \textbf{Finestra.} Tutti gli spostamenti sono vicini a zero e dentro le barre ($|\Delta|\le0.0014$): gruppo $0.0014$ per entrambi.
\item \textbf{Fondo.} Forma condivisa tra cariche ($-0.0037$ e $-0.0040$) e simmetrico $A_b=0$ ($-0.0028$ e $-0.0031$); la singola esponenziale ($-0.0037$, $-0.0039$) è disegnata vuota ed esclusa. Gruppo $0.0037$ (bayesiano) e $0.0040$ (frequentista).
\item \textbf{Segnale.} Crystal Ball ($-0.0002$, $-0.0006$) e $\mu,\sigma$ per carica ($\approx0$): gruppo $0.0002$ e $0.0006$.
\item \textbf{Preselezione.} $\mathtt{ProbK}>0.6$ dà $+0.0058$ con $\sigma_B\approx0.0050$ ($1.2\,\sigma_B$): il punto più lontano dallo zero, ma la sua barra contiene lo zero a poco più di $1\sigma_B$. Gruppo $0.0058$.
\item \textbf{Totale.} $\sqrt{0.0014^2+0.0037^2+0.0002^2+0.0058^2}=0.0070$ (bayesiano) e $0.0072$ (frequentista).
\end{enumerate}

\sez{Limiti e orale}
\begin{itemize}
\item La sistematica dominante (preselezione) è anche la più compatibile con una fluttuazione: l'incertezza sistematica è \emph{conservativa}. La Fig.~XX lo mostra con la scansione.
\item \textbf{``Perché il massimo per gruppo e non la somma?''} Le varianti di un gruppo sono alternative per la stessa scelta; sommarle conterebbe più volte lo stesso effetto.
\item \textbf{``Perché in quadratura?''} Gruppi diversi sono scelte indipendenti: le incertezze indipendenti si sommano in quadratura.
\item \textbf{``Come si distingue sistematico da statistico?''} Con $\sigma_B$ per i campioni annidati; per le varianti sullo stesso campione ogni spostamento è per definizione sistematico (il campione è identico).
\item Le $13$ varianti sono una scelta dell'analista: l'incertezza sistematica è l'incertezza \emph{dato questo insieme}, non un limite assoluto.
\end{itemize}

% ======================================================================== NEW syst_pid_scan
\figh{Fig.\ XX --- Scansione dei tagli di identificazione (App.\ E.6)}{syst_pid_scan}{0.30}
\sez{Il problema}
La preselezione domina la sistematica con una sola variante ($\mathtt{ProbK}>0.6$, $+0.0058$). Quel valore è un effetto vero del taglio o una fluttuazione statistica del sottocampione? Una singola variante non lo dice: serve guardare l'andamento \emph{lungo il taglio}.

\sez{Teoria}
I tagli $\mathtt{ProbK}$ (lower) e $\mathtt{ProbPi}$ (upper) sono definiti dalla verosimiglianza di identificazione di particella (PID). Stringere $\mathtt{ProbK}$ aumenta la purezza in kaoni e riduce il campione (campioni annidati); allentarlo lo allarga e fa entrare candidati con identificazione meno sicura, in particolare decadimenti mal identificati come $B^\pm\to K^\pm\pi^+\pi^-$ e $\pi^\pm K^+K^-$, che non sono nel modello di massa e possono avere un'asimmetria diversa.

Ci sono due comportamenti attesi. (i) Se lo spostamento è una \emph{fluttuazione}, al variare del taglio i punti scattano attorno allo zero con ampiezza $\sim\sigma_B$, senza andamento. (ii) Se c'è un \emph{effetto sistematico} dovuto a un fondo mal identificato, lo spostamento cresce in modo monotòno man mano che il taglio si allenta (più contaminazione).

\sez{Come si legge}
\begin{enumerate}
\item \textbf{(a)} Asse $x$: valore del taglio inferiore su $\mathtt{ProbK}$ ($0.30$--$0.85$); \textbf{(b)} taglio superiore su $\mathtt{ProbPi}$ ($0.15$--$0.70$). Asse $y$ comune: spostamento di $\Araw$ (fit frequentista) rispetto al nominale.
\item Linea verde con punti: lo spostamento ad ogni taglio. Linea verticale punto-tratto: taglio nominale ($0.5$), dove lo spostamento è per definizione $0$.
\item Fasce grigie: $\pm\sigma_B$ (scura) e $\pm2\sigma_B$ (chiara), l'ampiezza attesa dalla sola statistica del campione annidato; si annullano al nominale e si allargano stringendo il taglio (meno candidati).
\item Cosa guardare: i punti stanno dentro le fasce? C'è un andamento sistematico da un lato?
\end{enumerate}

\sez{Cosa dice}
\begin{enumerate}
\item \textbf{(a) Lato stretto ($\mathtt{ProbK}>0.5$).} I punti scattano attorno a zero senza andamento: $+0.0009$, $+0.0058$ ($1.2\sigma_B$), $-0.0017$, $-0.0112$ ($-1.3\sigma_B$), $-0.0054$, $+0.0010$, $+0.0022$: segni alternati, tutti entro $1.3\sigma_B$. Il valore $+0.0058$ della variante $\mathtt{ProbK}>0.6$ \emph{non} è un trend: il vicino $0.65$ ha segno opposto.
\item \textbf{(a) Lato largo ($\mathtt{ProbK}<0.5$).} Lo spostamento cresce in modo monotòno: $-0.0019$, $-0.0032$, $-0.0067$, $-0.0111$ per $\mathtt{ProbK}>0.45$, $0.40$, $0.35$, $0.30$, l'ultimo a $2.5\sigma_B$ (fuori dalla fascia $2\sigma_B$). È l'andamento atteso per contaminazione da decadimenti mal identificati.
\item \textbf{(b) $\mathtt{ProbPi}$.} Piatto: tutti gli spostamenti tra $-0.001$ e $+0.004$, dentro $\pm1.3\sigma_B$.
\item \textbf{Conseguenza.} La sistematica dominante è consistente con una fluttuazione di un campione annidato sul lato stretto; sul lato largo il modello vede un effetto reale, che \emph{sostiene} il taglio nominale (tenerlo a $0.5$ evita la regione dove lo spostamento è monotòno).
\end{enumerate}

\sez{Limiti e orale}
\begin{itemize}
\item I punti sono \emph{correlati} (campioni annidati): non sono dodici misure indipendenti.
\item \textbf{``Allora la sistematica dovuta al PID è zero?''} No: se fosse una pura fluttuazione, includerla è conservativo. Si include comunque perché $|\Delta|>\sigma_B$ non è escluso e perché il lato largo mostra un effetto vero.
\item \textbf{``Perché non un taglio più stretto?''} Si perde statistica ($\sigma$ cresce da $0.0085$ a $0.0099$ a $\mathtt{ProbK}>0.6$) senza guadagno dimostrabile.
\item Un'analisi con un modello del fondo mal identificato (da simulazione) potrebbe usare un campione più largo: è tra gli sviluppi futuri.
\end{itemize}

% ======================================================================== XVII
\figh{Fig.\ XXI --- Validazione della regola per campioni annidati (App.\ E.8)}{validation_barlow}{0.28}
\sez{Il problema}
Una variante di selezione (es.\ $\mathtt{ProbK}>0.6$) \emph{rimuove} candidati e sposta $\hat A$. Quanto di questo spostamento è un \emph{effetto sistematico} del modello e quanto è una fluttuazione statistica del sottocampione? Senza distinguerli, ogni variante conterebbe come sistematica e si sovrastimerebbe l'incertezza.

\sez{Teoria}
Sia $\hat A$ la stima sul campione nominale (efficiente) e $\hat A'$ quella su un suo sottocampione. Per uno stimatore efficiente vale $\mathrm{Cov}(\hat A',\hat A)=\mathrm{Var}(\hat A)$ (qualunque altro stimatore non polarizzato ha covarianza con quello di minima varianza pari alla varianza di quest'ultimo). Allora
\[
\mathrm{Var}(\hat A'-\hat A)=\sigma'^2+\sigma^2-2\sigma^2=\sigma'^2-\sigma^2\equiv\sigma_B^2.
\]
Una variante conta solo se $|\Delta|>\sigma_B$ (Barlow). Si applica ai cambi di campione (finestra, preselezione); non alle varianti dello stesso campione (forma del fondo, segnale), dove si usa lo spostamento intero.

\sez{Come si legge}
\begin{enumerate}
\item Pannelli (a) $\mathtt{ProbK}>0.6$ e (b) $\mathtt{ProbPi}<0.4$: istogramma dello spostamento $\hat A'-\hat A$ nelle $300$ repliche bootstrap, diviso per $\sigma_B$ mediano.
\item La curva rossa è una normale standard di \emph{riferimento}; in alto: deviazione standard dello spostamento e $\sigma_B$ mediano.
\item \textbf{Cosa si controlla è la larghezza:} se l'istogramma è largo $\approx1$ in unità di $\sigma_B$ (o, equivalentemente, se $\mathrm{sd}(\text{shift})\approx\sigma_B$) la regola è corretta.
\item La \emph{posizione} dell'istogramma non è zero: è lo spostamento osservato nei dati, che il bootstrap preserva.
\end{enumerate}

\sez{Cosa dice}
\begin{enumerate}
\item \textbf{Larghezza.} $\mathrm{sd}(\text{shift})=0.0051$ e $0.0025$ contro $\sigma_B$ mediano $0.0050$ e $0.0023$: rapporti $1.02$ e $1.09$. La formula $\sigma_B=\sqrt{\sigma'^2-\sigma^2}$ descrive la fluttuazione statistica osservata nei bootstrap.
\item \textbf{Posizione.} L'istogramma è centrato a $+1.2$ e $+0.9$ in unità di $\sigma_B$: gli spostamenti medi sono $0.0060$ e $0.0020$, quelli osservati nei dati.
\item \textbf{Valutazione.} $\mathtt{ProbK}>0.6$ ha shift $+0.0058$ e $\sigma_B=0.0050$: $1.2\sigma_B$, compatibile con una fluttuazione. La variante non è un effetto sistematico \emph{provato}, ma viene comunque inclusa nell'incertezza (conservativo).
\item \textbf{Conseguenza.} La regola è corretta: la sistematica dominante (preselezione, $0.0058$) è conservativa, non gonfiata da fluttuazioni statistiche non riconosciute.
\end{enumerate}

\sez{Limiti e orale}
\begin{itemize}
\item L'istogramma \emph{non} è centrato a zero: la normale a zero è solo un riferimento di larghezza; il controllo è sulla larghezza, non sulla posizione (da dire nella didascalia).
\item \textbf{``Perché sottrarre in quadratura $\sigma'^2-\sigma^2$ e non sommare?''} Perché i due campioni sono correlati (uno contiene l'altro): la differenza ha varianza minore della somma.
\item \textbf{``La regola vale per qualunque stimatore?''} Solo per stimatori efficienti (ML asintotico); per stimatori non efficienti $\mathrm{Cov}\neq\mathrm{Var}$ e la formula non è esatta.
\item \textbf{``Come combini le varianti?''} Per gruppo si prende la massima $|\Delta|$ sopra $\sigma_B$ e si sommano i gruppi in quadratura: $0.0070$ (bayesiano), $0.0072$ (frequentista).
\end{itemize}

% ======================================================================== XVIII
\figh{Fig.\ XXII --- Validità del teorema di Wilks (App.\ F)}{validation_null_wilks}{0.28}
\sez{Il problema}
Il $p$-value $2\times10^{-7}$ dipende dall'assunzione $q(0)\sim\chi^2_1$ fino a $q=27$. Si può verificare questa assunzione, e fin dove?

\sez{Teoria}
\emph{Funzione di sopravvivenza:} $P(q\ge q_{\rm obs})=1-F(q_{\rm obs})$. \emph{Trasformazione integrale di probabilità:} se $q\sim F$ allora $p=1-F(q)\sim U(0,1)$; quindi, se la distribuzione è quella giusta, il grafico quantile--quantile dei $p$-value contro quantili uniformi sta sulla diagonale. \emph{Test di Kolmogorov--Smirnov}: massima distanza tra funzione di ripartizione empirica e teorica.

\emph{Condizioni di Wilks:} (i)~$A=0$ è interno al dominio $(-1,1)$ (nessun bordo, a differenza di un yield $\ge0$); (ii)~ipotesi annidate; (iii)~$N$ grande; (iv)~modello identificabile sotto $H_0$.

\sez{Come si legge}
\begin{enumerate}
\item \textbf{(a)} Asse $x$: $q(0)$ (scala log, da $10^{-3}$ a $10^{1}$); asse $y$: $P(q\ge q_{\rm obs})$ (scala log, da $10^{-3}$ a $1$). Curva verde acqua = $293$ toy a $A_{\rm true}=0$; curva rossa = $\chi^2_1$ (Wilks). Devono sovrapporsi.
\item \textbf{(b)} Asse $x$: quantili uniformi; asse $y$: $p$-value di $t(A_{\rm true})$ ordinati. Sulla diagonale se la distribuzione è corretta. Due curve: valore fittato ($n=985$) e tutti i valori iniettati.
\item \textbf{(c)} Istogramma di $t(A_{\rm true})$ al valore fittato contro $\chi^2_1$; la linea punteggiata a $t=1$ è il limite al $68\%$.
\item Cosa guardare: deviazioni sistematiche dalla diagonale in (b) (curve sopra o sotto), o code diverse in (a).
\end{enumerate}

\sez{Cosa dice}
\begin{enumerate}
\item \textbf{(a)} Le due curve si sovrappongono fino a $q\approx5$; la coda dei toy si interrompe a $q\approx10$ (con $293$ toy il livello minimo risolvibile è $1/293\approx3\times10^{-3}$). Kolmogorov--Smirnov $p=0.79$.
\item \textbf{(b)} La curva sta sulla diagonale per tutto l'intervallo: i $p$-value sono uniformi, KS $p=0.96$ su tutti i valori iniettati.
\item \textbf{(c)} L'istogramma segue $\chi^2_1$, compreso il picco a $0$.
\item \textbf{Code.} $2$ toy su $293$ oltre $q=9$ contro $0.8$ attesi: compatibile.
\item \textbf{Conseguenza.} $q(0)$ è distribuita come $\chi^2_1$ nel regime verificabile ($q\lesssim10$), quindi taglia e copertura sono giuste; il valore osservato $q=27$ è un'estrapolazione \emph{plausibile} ma non verificata direttamente.
\end{enumerate}

\sez{Limiti e orale}
\begin{itemize}
\item \textbf{``Come sai che $5.2\sigma$ è corretto se i toy arrivano solo a $3\sigma$?''} Non lo si può verificare direttamente (servirebbero $\sim5\times10^6$ toy per $p=2\times10^{-7}$); si estrapola la forma $\chi^2_1$, che è giustificata dal teorema e confermata fin dove si può.
\item Un test di Kolmogorov--Smirnov ha poca potenza nelle code: per questo si guarda anche il grafico (a) in scala log.
\item \textbf{``Wilks vale se il parametro è al bordo?''} No: la distribuzione diventa un misto di $\chi^2_0$ e $\chi^2_1$ (Chernoff); non è il caso di $A$.
\item \textbf{``Differenza tra (a) e (b)?''} (a) guarda le code con il valore osservato a $A_{\rm true}=0$; (b) verifica l'intera distribuzione per tutti gli $A_{\rm true}$ iniettati.
\end{itemize}

% ======================================================================== XIX
\figh{Fig.\ XXIII --- Iniezione di asimmetrie: bias, efficienza, potenza (App.\ F)}{validation_injection_power}{0.26}
\sez{Il problema}
Lo stimatore $\hat A$ è corretto, e l'incertezza di Hesse è quella vera, \emph{per ogni} valore dell'asimmetria iniettata, non solo per quello trovato nei dati? Se il bias dipendesse da $A$, il risultato sarebbe difficile da fidare.

\sez{Teoria}
\emph{Cramér--Rao:} per uno stimatore non polarizzato $\mathrm{Var}(\hat\theta)\ge1/I(\theta)$ con $I=-\E[\partial^2\ln L/\partial\theta^2]$ l'informazione di Fisher. Lo stimatore di massima verosimiglianza è asintoticamente non polarizzato, normale e \emph{efficiente} (raggiunge il limite). L'incertezza di Hesse è $[-\partial^2\ln L/\partial A^2]^{-1/2}$ valutata al massimo. Se la dispersione dei toy diviso la $\sigma_{\rm Hesse}$ media è $\approx1$, lo stimatore è efficiente e l'incertezza è corretta; il bias ha incertezza standard $\sigma/\sqrt n=0.0085/\sqrt{293}=0.0005$.

\sez{Come si legge}
\begin{enumerate}
\item Asse $x$ comune: $A_{\rm true}$ iniettato ($-0.1,-0.0441,-0.02,0,+0.05$).
\item \textbf{(a)} $y=\langle\hat A\rangle-A_{\rm true}$ con errore standard: i punti devono stare su $0$ (linea punteggiata).
\item \textbf{(b)} $y$ = due rapporti: $\mathrm{sd}(\hat A)/\langle\sigma_{\rm Hesse}\rangle$ (quadrati) e deviazione standard dei pull (cerchi): devono stare su $1$.
\item \textbf{(c)} Probabilità di rigettare $A=0$ con $q>3.84$ e $q>9$: la stessa informazione della Fig.~VI.
\item Cosa guardare: andamenti sistematici con $A_{\rm true}$ (una deriva indicherebbe un bias dipendente dal valore).
\end{enumerate}

\sez{Cosa dice}
\begin{enumerate}
\item \textbf{(a) Bias.} Al più $0.0008$ in valore assoluto ($1.6$ errori standard), tra $\approx-0.0008$ ($A_{\rm true}=-0.02$) e $+0.0002$ ($A_{\rm true}=0$): nessuna dipendenza sistematica, i punti scattano attorno a zero.
\item \textbf{(b) Efficienza.} Rapporti tra $0.94$ e $1.05$: la dispersione reale dello stimatore coincide (entro $6\%$) con l'incertezza di Hesse; la larghezza dei pull ha lo stesso andamento.
\item \textbf{(c)} Potenza vedi Fig.~VI.
\item \textbf{Conseguenza.} Lo stimatore ML è non polarizzato e raggiunge il limite di Cramér--Rao a questa statistica, e l'incertezza di Hesse ($0.0085$) è l'incertezza statistica corretta a ogni $A$ iniettato.
\end{enumerate}

\sez{Limiti e orale}
\begin{itemize}
\item L'ultimo punto ($+0.05$) ha copertura al $68\%$ pari a $0.61\pm0.03$, $2.5\sigma$ sotto: su cinque valori una deviazione di questa entità non è sorprendente (con $5$ test indipendenti, la probabilità di vederne una $\ge2.5\sigma$ è $\sim6\%$).
\item \textbf{``Cos'è l'efficienza di uno stimatore?''} Il rapporto tra la varianza minima (Cramér--Rao) e la sua varianza; $1$ = efficiente.
\item \textbf{``Perché è importante il bias dipende da $A$?''} Perché si misura $A$ vicino a $-0.044$, ma se fosse $\sim0$ o $\sim-0.1$ si vuole la stessa affidabilità (verifica che $\hat A$ non sia compresso verso lo $0$).
\end{itemize}

% ======================================================================== XX
\figh{Fig.\ XXIV --- $\chi^2$ binnato e binning (App.\ F)}{validation_gof_toys}{0.30}
\sez{Il problema}
Il $p$-value del $\chi^2$ è $5\times10^{-6}$ in bin da $10\MeV$ e $0.03$ in bin da $2.5\MeV$ per \emph{lo stesso} fit. È una proprietà del modello o del test? Senza saperlo non si può dire se ``il modello è inadeguato'' sia una conclusione robusta.

\sez{Teoria}
Il $\chi^2$ di Pearson segue $\chi^2_{k-p}$ se i conteggi attesi sono grandi ($\gtrsim5$) e i parametri sono stimati minimizzando lo stesso $\chi^2$ sui dati binnati. Se i parametri sono stimati con la verosimiglianza \emph{unbinned} (come qui), la distribuzione asintotica è compresa tra $\chi^2_{k-p}$ e $\chi^2_{k-1}$ (Chernoff--Lehmann): la differenza è piccola per $k\gg p$. Media $\approx k-p$, varianza $2(k-p)$; $\sqrt{2\cdot88}=13.3$.

Il \emph{potere} del test dipende dal binning: una discrepanza su scala $\Delta m$ è diluita in bin molto più stretti, dove le fluttuazioni di Poisson dominano; in bin molto larghi la discrepanza si media. Il $\chi^2$ non ha un binning ``giusto'': il $p$-value è proprietà della coppia (modello, binning).

\sez{Come si legge}
\begin{enumerate}
\item Asse $x$: valore del $\chi^2$ binnato; asse $y$: densità.
\item Istogramma azzurro = $\chi^2$ di $1000$ toy generati dal modello fittato e rifittati; curva rossa = $\chi^2$ con i gradi di libertà corrispondenti (88 in (a), 388 in (b)); linea rossa verticale = valore dei dati (160 e 442).
\item (a) $2\times50$ bin da $10\MeV$; (b) $2\times200$ bin da $2.5\MeV$.
\item Cosa guardare: (i) l'istogramma dei toy segue la curva rossa? (ii) dove cade la linea dei dati rispetto al corpo dell'istogramma?
\end{enumerate}

\sez{Cosa dice}
\begin{enumerate}
\item \textbf{I toy seguono la $\chi^2$.} Media $87.6$, $\sigma=13.3$ con $88$ dof; $387.6$ e $27.1$ con $388$ dof: i $p$-value della $\chi^2$ sono accurati per questo modello e questa statistica.
\item \textbf{(a)} Dati a $159.5$: $5.4$ deviazioni standard sopra i toy ($p=4.8\times10^{-6}$), fuori dalla coda: nessun toy lo raggiunge.
\item \textbf{(b)} Dati a $441.7$: $2.0$ deviazioni standard ($p=0.031$, nei toy $0.032$): dentro la coda destra dell'istogramma.
\item \textbf{Perché.} Il misfit ha scala $\sim10\MeV$ (i fianchi del picco): in bin da $10\MeV$ si accumula in pochi bin; in bin da $2.5\MeV$ si diluisce su molti bin con fluttuazioni di Poisson.
\item \textbf{Conseguenza.} Il $p$-value del $\chi^2$ è proprietà del binning quanto del modello; la conclusione ``il modello non è perfetto'' resta, ma con quantificazioni diverse: $5.4\sigma$ o $2.0\sigma$.
\end{enumerate}

\sez{Limiti e orale}
\begin{itemize}
\item ``Il modello non è adeguato'' e ``il modello è accettabile'' sono entrambi sostenibili a seconda del binning; lo stesso vale per il Bayes factor con la larghezza del prior.
\item \textbf{``Qual è il binning giusto?''} Non esiste; si sceglie in anticipo (qui $10\MeV$, comparabile con la risoluzione $\sigma=18\MeV$) e si riportano più scelte.
\item \textbf{``Alternative al $\chi^2$ binnato?''} Test unbinned (Kolmogorov--Smirnov, Anderson--Darling) o il PPC (Fig.~XV).
\item \textbf{``Perché $\chi^2/{\rm ndof}=1.86$ non invalida $\Araw$?''} Perché il misfit è sulle code, comuni alle due cariche (Fig.~III) e $\Araw$ è insensibile alla loro forma.
\end{itemize}


% ======================================================================== NEW dalitz_perm
\figh{Fig.\ XXV --- Test di permutazione delle cariche nel Dalitz (App.\ G.4)}{dalitz_permutation}{0.34}
\sez{Il problema}
Il test di omogeneità $\chi^2=\sum_iS_i^2=80.8$ per $31$ gradi di libertà (Fig.~VII) assume che ogni $S_i$ sia gaussiano standard. Quella assunzione vale per bin con $\ge20$ candidati ma non è esatta. Esiste una verifica \emph{senza} assunzioni distribuzionali?

\sez{Teoria}
\emph{Test di permutazione.} Si costruisce la distribuzione della statistica sotto l'ipotesi nulla ``l'asimmetria è uniforme'' riassegnando a caso le cariche $q_i$ ai candidati, mantenendo fissi la posizione nel Dalitz di ciascuno (e quindi l'occupazione dei bin) e il numero totale di $B^+$ e $B^-$. Sotto $H_0$ la carica è \emph{scambiabile} tra candidati, quindi ogni permutazione è un'estrazione equiprobabile dalla distribuzione nulla. Si ripete $n_{\rm perm}$ volte e il $p$-value è
\[
p=\frac{1+\#\{\chi^2_{\rm perm}\ge\chi^2_{\rm obs}\}}{1+n_{\rm perm}},
\]
con il minimo risolvibile $1/(1+n_{\rm perm})\approx5\times10^{-4}$ per $2000$ permutazioni. È \emph{esatto} (non asintotico) e non richiede ipotesi gaussiane.

\sez{Come si legge}
\begin{enumerate}
\item Asse $x$: valore della statistica $\chi^2=\sum_iS_i^2$ sulla griglia $8\times8$ ($32$ bin popolati); asse $y$: densità.
\item Istogramma azzurro: la statistica per $2000$ permutazioni casuali delle cariche. Curva nera tratteggiata: $\chi^2$ con $31$ gradi di libertà. Linea rossa: valore osservato ($80.8$).
\item Cosa guardare: (i) l'istogramma segue la curva? (conferma che $S_i$ è ben normalizzato); (ii) dove cade la linea rossa rispetto alla coda destra?
\end{enumerate}

\sez{Cosa dice}
\begin{enumerate}
\item \textbf{Normalizzazione corretta.} L'istogramma segue la $\chi^2_{31}$ (media $\approx31$, deviazione $\approx8$): i $p$-value gaussiani del test sono affidabili, e il test è già corretto per il numero di bin.
\item \textbf{Il dato è fuori.} Nessuna delle $2000$ permutazioni raggiunge $80.8$ (il massimo è $\approx72$): $p<5\times10^{-4}$, in accordo con $p=3\times10^{-6}$ del $\chi^2$ asintotico (che l'esperimento non può risolvere oltre $1/2000$).
\item \textbf{Conseguenza.} La non-uniformità dell'asimmetria locale è un fatto dei dati, non dipende dall'approssimazione gaussiana. Il valore osservato è $\approx6$ deviazioni standard sopra la media nulla.
\end{enumerate}

\sez{Limiti e orale}
\begin{itemize}
\item La permutazione \emph{preserva} l'occupazione dei bin ma rompe ogni correlazione tra carica e posizione: sotto $H_0$ l'asimmetria è globale e uniforme. Non testa asimmetrie \emph{strumentali} che dipendono dalla posizione (detector, impulso): per quello c'è la divisione per polarità.
\item \textbf{``Perché non solo il $\chi^2$ asintotico?''} Perché le condizioni gaussiane sono approssimate; la permutazione le verifica.
\item \textbf{``Cos'è la scambiabilità?''} L'ipotesi che la distribuzione congiunta dei dati non cambi permutando le etichette: è ciò che $H_0$ afferma.
\item \textbf{``Un valore $p<1/n_{\rm perm}$?''} Si riporta solo come limite superiore: $p<5\times10^{-4}$.
\item Il test sulla simulazione phase-space ($\chi^2=40.0/30$, $p=0.11$) verifica che la procedura non inventi asimmetria dove non c'è.
\end{itemize}

% ======================================================================== XXI

\figh{Fig.\ XXVI --- Asimmetria nella banda del $\phi(1020)$ e nei bin significativi (App.\ G.5)}{dalitz_phi_band}{0.30}

\sez{Il problema}
Il corpo afferma che l'asimmetria negativa a basso $m^2_{\rm lo}$ \emph{non} è prodotta dal $\phi(1020)$, la risonanza stretta a $m^2_{\rm lo}\approx1.04\GeV^2$ che sta proprio accanto alla regione. La verifica è una selezione: si isolano i candidati del $\phi$ e si misura l'asimmetria con e senza.

\sez{Teoria}
Il $\phi(1020)$ ($m=1019.46\MeV$, larghezza $4.25\MeV$) decade in $K^+K^-$; nel Dalitz appare come banda stretta a $m^2_{\rm lo}=m_\phi^2$. In $B^\pm\to K^\pm\phi$ non c'è interferenza tra ampiezze con fase forte diversa che produca una grande asimmetria diretta, quindi ci si attende un'asimmetria piccola. Se l'effetto fosse dovuto al $\phi$, rimuovendo la sua banda l'asimmetria della regione dovrebbe ridursi; se non lo è, resta o aumenta.

L'asimmetria di conteggio in un insieme di $n$ candidati è $A=(N^--N^+)/n$ con $\sigma=\sqrt{(1-A^2)/n}$.

\sez{Come si legge}
\begin{enumerate}
\item \textbf{(a)} Istogramma di $m_{\rm lo}(K^+K^-)$ in MeV (bin da $5\MeV$): il picco del $\phi$ a $1019.5$ (linea tratteggiata); la fascia grigia è la banda usata, $|m-m_\phi|<10\MeV$.
\item \textbf{(b)} Per ciascun insieme di candidati (righe): punto con barra = asimmetria di conteggio e incertezza binomiale; a destra il numero di candidati. La linea punteggiata è il valore globale ($-0.0432$), la continua è $0$.
\item Cosa guardare: l'insieme ``banda del $\phi$'' è a zero? Le regioni senza $\phi$ restano negative?
\end{enumerate}

\sez{Cosa dice}
\begin{enumerate}
\item \textbf{Banda del $\phi$} ($1964$ candidati): $-0.004\pm0.023$, compatibile con zero.
\item \textbf{Appena sotto, fuori dalla banda} ($m^2_{\rm lo}<2\GeV^2$, $4833$ candidati): $-0.120\pm0.014$: l'asimmetria è nella regione, non nel picco del $\phi$.
\item \textbf{Bin significativi} ($3091$ candidati): $-0.128\pm0.018$; \emph{senza} i candidati del $\phi$ ($2738$): $-0.145\pm0.019$: l'asimmetria \emph{aumenta} se si toglie il $\phi$; i soli $353$ candidati del $\phi$ in quella regione: $+0.003\pm0.053$.
\item \textbf{Conseguenza.} L'effetto non viene dal $\phi$; il $\phi$ lo \emph{diluisce}. Coerente con l'interferenza delle risonanze a basso $m^2_{\rm lo}$.
\end{enumerate}

\sez{Limiti e orale}
\begin{itemize}
\item Sono conteggi: il fondo (più simmetrico) diluisce tutti i valori; nel $\phi$ il fondo è piccolo ma non nullo.
\item La banda è $\pm10\MeV$ attorno alla massa nominale: una banda più larga o stretta cambia di poco i numeri (non variata nel report).
\item \textbf{``Perché non hai fittato anche la banda del $\phi$?''} Con $1964$ candidati il fit è possibile; il conteggio basta a mostrare l'assenza dell'effetto.
\end{itemize}

\figh{Fig.\ XXVII --- Regioni di confidenza e credibili in $(A,A_b)$ (App.\ H.4)}{validation_region_Ab}{0.40}
\sez{Il problema}
Finora il confronto freq/Bayes è stato in una dimensione. Per due parametri d'interesse insieme, $A$ e $A_b$ (le due asimmetrie, di segnale e di fondo), i due metodi danno ancora le stesse regioni? E dove sta $A_b=0$ rispetto al risultato?

\sez{Teoria}
Per $k$ parametri d'interesse la statistica di Wilks segue $\chi^2_k$. Le soglie per il contenuto $\mathrm{CL}$ sono i quantili di $\chi^2_k$: per $k=2$, $-2\ln(1-0.6827)=2.30$ e $-2\ln(1-0.95)=5.99$ (a differenza di $1$ e $3.84$ in 1D). \emph{Regione HPD} (highest posterior density): la regione più piccola che contiene la probabilità desiderata, con densità posteriore sopra una soglia. Per una gaussiana bivariata sia la regione di profilo sia la HPD sono ellissi, inclinate da $\rho$.

Attenzione: la \emph{proiezione} su $A$ della regione bidimensionale al $68\%$ è più larga dell'intervallo monodimensionale al $68\%$ (la prima corrisponde a $q\le2.30$, la seconda a $q\le1$).

\sez{Come si legge}
\begin{enumerate}
\item Assi: $x$ = $\Araw$ (da $-0.075$ a $0.005$), $y$ = $A_b$ (da $-0.11$ a $0.03$).
\item Due famiglie di contorni: frazione \textbf{verde acqua} = frequentista (profilo con soglie $\chi^2_2$); \textbf{viola} = bayesiana (HPD sulla catena). Continuo = $68\%$, tratteggiato = $95\%$.
\item Il punto centrale è il best-fit / mediana.
\item Le due linee punteggiate sono $A=0$ (verticale, a destra) e $A_b=0$ (orizzontale, in alto).
\item Cosa guardare: i contorni verdi e viola coincidono? Le linee punteggiate cadono dentro o fuori le regioni?
\end{enumerate}

\sez{Cosa dice}
\begin{enumerate}
\item \textbf{Coincidenza.} Le ellissi frequentista e bayesiana si sovrappongono quasi perfettamente a $68\%$ e a $95\%$: i due metodi danno la stessa regione bidimensionale.
\item \textbf{Forma.} Quasi circolari con una leggera inclinazione: correlazione $\rho=-0.16$ (verosimiglianza), $-0.17$ (posterior), $-0.17$ (bootstrap).
\item \textbf{$A=0$.} Sta a $\approx0.04$ dal bordo destro del $95\%$, cioè oltre $5\sigma$: fuori da entrambe le regioni.
\item \textbf{$A_b=0$.} Cade sul \emph{bordo} della regione al $95\%$: $A_b=-0.0402\pm0.0176$ ($2.3\sigma$ in 1D), consistente con il Bayes factor che lo trova marginale.
\item \textbf{Conseguenza.} L'asimmetria del segnale è decisamente diversa da zero; quella del fondo è un indizio debole. I due metodi concordano anche in due dimensioni.
\end{enumerate}

\sez{Limiti e orale}
\begin{itemize}
\item Qui $A_b$ è un parametro d'interesse; nelle altre figure è un nuisance (profilato o marginalizzato). Cambiare ``cosa è d'interesse'' cambia le soglie ma non l'ellisse.
\item \textbf{``Perché $2.30$ e non $1$?''} Perché a due parametri la regione al $68\%$ richiede $q\le2.30$: la $\chi^2_2$ ha cumulativa $1-e^{-q/2}$.
\item \textbf{``HPD o equal-tailed in 2D?''} In 2D non c'è un'ovvia coda: si usa la regione di densità più alta.
\item \textbf{``Hanno la stessa interpretazione?''} No: la frequentista ha copertura, la bayesiana è una probabilità sui parametri.
\end{itemize}

% ======================================================================== XXII
\figh{Fig.\ XXVIII --- I due metodi a bassa statistica (App.\ H.11)}{bayes_nominal_lowstat}{0.30}
\sez{Il problema}
Con $N\approx2\times10^4$ i due metodi coincidono per costruzione (Bernstein--von Mises). Esiste un regime, di pochi eventi, in cui divergono? Cercarlo mostra \emph{dove} è la vera differenza tra i due approcci.

\sez{Teoria}
I due metodi differiscono quando la verosimiglianza è molto non gaussiana (pochi eventi, parametri al bordo) o il prior è informativo. $\hat A=(N_s^--N_s^+)/(N_s^-+N_s^+)$ è un rapporto di conteggi di Poisson: per yield $\gtrsim10$--$20$ la verosimiglianza in $A$ è già quasi gaussiana, quindi Hesse $\approx$ Minos $\approx$ credibile. L'errore binomiale sulla copertura con $300$ toy è $\sqrt{0.68\cdot0.32/300}=0.027$.

\sez{Come si legge}
\begin{enumerate}
\item Asse $x$ comune (scala log): numero di candidati per dataset, $N=50,150,500,2000$. Asse $y$: copertura osservata.
\item \textbf{(a)} Livello $68\%$; \textbf{(b)} livello $95\%$. Linea orizzontale tratteggiata = nominale ($0.683$ e $0.954$).
\item Tre punti per ogni $N$ (leggermente sfalsati): credibile bayesiano (viola, cerchio), profilo/Minos (verde acqua, quadrato), Hesse (grigio, triangolo), con barre binomiali.
\item Cosa guardare: i punti sono compatibili col nominale entro le barre? I tre metodi coincidono tra loro?
\end{enumerate}

\sez{Cosa dice}
\begin{enumerate}
\item \textbf{Accordo col nominale.} Tutti i punti sono compatibili col nominale entro $\sim1\sigma$, anche a $N=50$ (una decina di candidati per carica).
\item \textbf{Accordo tra metodi.} Credibile, Minos e Hesse danno coperture quasi uguali a ogni $N$; la differenza è sotto le barre.
\item \textbf{$N=2000$.} Tutti sovra-coprono ($\approx0.72$ al $68\%$, $0.98$--$0.99$ al $95\%$): effetto della statistica finita dell'insieme di repliche, \emph{uguale} per i tre metodi, non di uno solo.
\item \textbf{Mediana contro MLE.} Differiscono mediamente del $5\%$ dell'incertezza statistica a $N=50$ e meno sopra.
\item \textbf{Conseguenza.} Nemmeno a $N=50$ i due metodi divergono: i numeri sono gli stessi, ciò che cambia è il significato dell'intervallo e il test d'ipotesi.
\end{enumerate}

\sez{Limiti e orale}
\begin{itemize}
\item Usa il modello minimale a sei parametri (gaussiana più una esponenziale, App.~B), identificabile a bassa statistica; a $N=50$ la catena è lunga $\approx10\tau$, quindi l'errore Monte Carlo sugli estremi è maggiore.
\item \textbf{``Quando divergerebbero?''} Con parametro al bordo (yield $\to0$), poche decine di eventi totali, prior informativo, o verosimiglianza multimodale.
\item \textbf{``Perché un modello più semplice?''} Il nominale a $12$ parametri non è identificabile con $50$ candidati.
\item I due metodi differiscono nel \emph{significato} e nei test d'ipotesi, non nei numeri: è il messaggio centrale del report.
\end{itemize}

\clearpage
\section*{Parte ML (App.\ I)}
Nelle figure seguenti i modelli sono addestrati con \emph{weak labels}: segnale = finestra $\pm2\sigma$ intorno al picco ($15\,130$ candidati), fondo = bande laterali ($1940$). ROC, AUC e calibrazione misurano quanto si separano queste due regioni di massa, non la purezza vera S/B. Il ML ha il ruolo di \emph{controllo di coerenza}: se la selezione cambiasse $\Araw$ oltre le fluttuazioni statistiche, l'asimmetria potrebbe essere un artefatto della selezione per PID e massa.

% ======================================================================== XXIII

\figh{Fig.\ XXIX --- Weak labels sullo spettro di massa (App.\ I.4)}{ml_weak_labels}{0.34}

\sez{Il problema}
Il classificatore ML non ha etichette vere. Come si definisce, e su quali candidati è addestrato? La figura mostra le regioni di massa da cui vengono le etichette.

\sez{Teoria}
\emph{Weak supervision / CWoLa.} Si dispone di due campioni \emph{misti} con frazioni di segnale diverse: la finestra di picco (alta frazione di segnale) e le bande laterali (bassa). Un classificatore che separa i due campioni è monotòno nel rapporto di verosimiglianza S/B (Fig.~XXX), purché le feature non portino informazione sulla massa e il fondo nelle bande sia rappresentativo di quello sotto il picco. Le etichette sono per definizione \emph{rumorose} (il picco contiene $\approx4\%$ di fondo, le bande qualche componente non combinatoria).

La zona $5100$--$5150\MeV$ è esclusa dalla classe di fondo perché lì sale il fondo parzialmente ricostruito, che ha cinematica diversa da quella del combinatorio sotto il picco: includerla violerebbe la condizione di rappresentatività.

\sez{Come si legge}
\begin{enumerate}
\item Asse $x$: $m(KKK)$ ($5100$--$5600\MeV$); asse $y$: candidati per $5\MeV$ (scala log).
\item Istogramma nero = tutti i candidati; blu = signal-like ($|m-\mu|<2\sigma$, $15\,130$); arancione = background-like (bande oltre $4\sigma$ e sopra $5150\MeV$, $1940$).
\item Grigio chiaro = intervallo non etichettato ($2$--$4\sigma$); tratteggiato = zona esclusa ($<5150\MeV$).
\item Cosa guardare: l'estensione relativa delle classi e come il fondo sotto il picco è campionato dalle bande.
\end{enumerate}

\sez{Cosa dice}
\begin{enumerate}
\item \textbf{Sbilanciamento.} $1940$ background-like contro $15\,130$ signal-like (rapporto $0.128$): da qui i pesi di classe $4.40$ e $0.56$.
\item \textbf{Lacuna.} L'intervallo $2$--$4\sigma$ non è etichettato per ridurre la contaminazione tra le classi.
\item \textbf{Conseguenza.} Il classificatore impara a distinguere \emph{picco} da \emph{bande}; il suo valore come discriminante S/B dipende dalle due condizioni sopra, verificate in Fig.~XXXIV (no sculpting) e discusse in Fig.~XLI (fondo piccato non separabile).
\end{enumerate}

\sez{Limiti e orale}
\begin{itemize}
\item Le etichette dipendono da $\mu,\sigma$ del fit bayesiano: sono propagate come costanti.
\item \textbf{``Perché non si usa la simulazione per le etichette?''} Il record contiene solo una simulazione phase-space del segnale senza fondo; per il fondo servono i dati.
\item \textbf{``Cosa succede ai $B$ parzialmente ricostruiti?''} Stanno sotto $5150\MeV$ (esclusi dalle etichette) ma il classificatore li promuove comunque (Fig.~XL).
\end{itemize}

\figh{Fig.\ XXX --- Curve ROC (App.\ I.5)}{ml_nominal_roc}{0.38}
\sez{Il problema}
Sono stati addestrati sei classificatori, da lineari ad ensemble e reti. Quale separa meglio le due regioni di massa, e la differenza è reale o dentro la fluttuazione del campione di test? Serve per scegliere il riferimento senza usare il test.

\sez{Teoria}
\emph{Classificazione.} Un classificatore stima $p(y=1\mid x)$ minimizzando una loss (entropia incrociata binaria, il negativo della log-verosimiglianza di Bernoulli). \emph{ROC:} per una soglia $t$ sullo score $s$, $\mathrm{TPR}(t)=P(s>t\mid\text{signal})$ e $\mathrm{FPR}(t)=P(s>t\mid\text{bkg})$; la ROC è la curva $(\mathrm{FPR},\mathrm{TPR})$ al variare di $t$. \emph{AUC} $=P(s_{\rm sig}>s_{\rm bkg})$ (statistica di Mann--Whitney): $0.5$ = casuale, $1$ = perfetto. La ROC è invariante per trasformazioni monotòne dello score, quindi non dipende dalla calibrazione.

Per Neyman--Pearson, il test più potente a efficienza data è il rapporto di verosimiglianza $p_s(x)/p_b(x)$: nessun classificatore può stare sopra la ROC del LR.

\emph{CWoLa} (classification without labels). Siano due campioni misti con frazioni di segnale $f_1> f_2$. Il classificatore ottimo tra $M_1$ e $M_2$ è
\[
\frac{p(M_1\mid x)}{p(M_2\mid x)}\propto\frac{f_1p_s(x)+(1-f_1)p_b(x)}{f_2p_s(x)+(1-f_2)p_b(x)},
\]
una funzione \emph{monotòna} di $p_s/p_b$: perciò ha la stessa ROC del discriminante S/B ottimo, a patto che $p_s$ e $p_b$ non dipendano dal campione (le feature non devono contenere la massa, Fig.~XL).

\emph{Modelli:} regressione logistica (confine lineare); albero (partizioni ortogonali, profondità $4$); Random Forest (\emph{bagging} di $300$ alberi, ognuno su bootstrap e su sottoinsiemi di feature: riduce la varianza); gradient boosting (alberi poco profondi aggiunti in sequenza lungo il gradiente della loss: riduce il bias); FFNN (due strati nascosti $32$, $16$, ReLU, dropout, Adam).

\sez{Come si legge}
\begin{enumerate}
\item Asse $x$: efficienza sul background-like (false positive rate), da $0$ a $1$; asse $y$: efficienza sul signal-like (true positive rate).
\item Ogni curva è un modello (colori in legenda con l'AUC); la diagonale tratteggiata è il classificatore casuale; il cerchio nero è il taglio $0.5$ del gradient boosting.
\item Più la curva sale verso l'angolo in alto a sinistra, meglio: alta efficienza sul segnale a bassa efficienza sul fondo.
\item Cosa guardare: (i) la posizione relativa delle curve; (ii) dove si sovrappongono; (iii) il punto di lavoro (cerchio).
\end{enumerate}

\sez{Cosa dice}
\begin{enumerate}
\item \textbf{Gruppi.} Le due regressioni logistiche (AUC $0.771$, sovrapposte) sono in basso; l'albero singolo ($0.824$) sta in mezzo; Random Forest ($0.860$), gradient boosting ($0.864$) e FFNN ($0.857$) sono sovrapposti in alto.
\item \textbf{Salto lineare $\to$ non lineare.} $+0.09$ di AUC: il confine tra le due regioni non è lineare nelle feature.
\item \textbf{Fra i non lineari.} Differenze $\le0.007$, sotto l'incertezza del test ($0.012$, bootstrap): non sono distinguibili. L'albero singolo è $0.04$ sotto: il \emph{bagging/boosting} riduce la varianza.
\item \textbf{Punto di lavoro.} Il cerchio a score $0.5$: tiene l'$88\%$ dei signal-like e il $28\%$ dei background-like.
\item \textbf{Conseguenza.} Il gradient boosting è il riferimento perché migliore in cross-validation ($0.866$ contro $0.863$ della RF), non sul test, per non usare il test per nessuna scelta.
\end{enumerate}

\sez{Limiti e orale}
\begin{itemize}
\item L'AUC misura la separazione delle etichette \emph{deboli}, non la purezza vera S/B. Con $15\,130$ contro $1940$ candidati il test contiene solo $388$ background-like: differenze sotto $\sim0.01$ non sono risolte.
\item \textbf{``Perché la ROC e non l'accuratezza?''} Perché è indipendente dalla soglia e dalle proporzioni di classe.
\item \textbf{``Come dai un'incertezza all'AUC?''} Bootstrap del test set ($\pm0.012$) e deviazione standard su $5$ fold in cross-validation ($\pm0.016$).
\item \textbf{``Perché non scegliere il modello sul test?''} Si introdurrebbe un ottimismo (winner's curse): la scelta è fatta in cross-validation.
\end{itemize}

% ======================================================================== XXIV
\figh{Fig.\ XXXI --- Importanza delle feature (App.\ I.5)}{ml_nominal_features}{0.50}
\sez{Il problema}
Il classificatore usa informazione fisica sensata (la topologia del decadimento: vertice spostato) o scorciatoie, in particolare la massa, che invaliderebbero il suo uso come selezione per l'asimmetria? Si guarda cosa pesa di più.

\sez{Teoria}
\emph{Impurità di Gini:} $G=1-\sum_cp_c^2$, con $p_c$ le frazioni di classe in un nodo. Uno split che divide il nodo in due figli riduce l'impurità di $\Delta G=G-\frac{n_L}{n}G_L-\frac{n_R}{n}G_R$. \emph{Importanza} di una feature nella RF = $\Delta G$ totale degli split su quella feature, pesato per il numero di campioni e mediato sugli alberi: misura l'\emph{uso} in split non lineari ed è \emph{condivisa} tra feature correlate.

\emph{Lasso:} regressione logistica con penalità $\alpha\|w\|_1$; la norma $L_1$ spinge i coefficienti a zero (selezione). Su feature standardizzate $|w_j|$ misura l'effetto di $x_j$ \emph{a parità delle altre}. Correlazione di rango di Spearman: correlazione di Pearson dei ranghi.

\emph{Variabili:} $\mathtt{IPChi2}$ = aumento di $\chi^2$ del vertice primario aggiungendo la traccia: grande per tracce che \emph{non} vengono dal PV, cioè dai figli del $B$ (vertice spostato). $\mathtt{B\_FlightDistance}$ = distanza PV--vertice del $B$: il $B^\pm$ ha $c\tau\approx490\,\mu$m e il boost di Lorentz lo porta a millimetri; combinazioni casuali di tracce no. $\mathtt{B\_VertexChi2}$ = qualità del fit del vertice del $B$ (bassa per $B$ veri).

\sez{Come si legge}
\begin{enumerate}
\item Due pannelli con lo stesso elenco di $20$ feature sull'asse $y$ (ordinate per importanza RF decrescente dall'alto).
\item \textbf{(a)} Lunghezza della barra = importanza di Gini della RF (somma $=1$).
\item \textbf{(b)} Lunghezza = $|\text{coefficiente Lasso}|$ della feature standardizzata. Dentro (b) in basso è riportata la correlazione di Spearman tra le due graduatorie.
\item Le barre in \emph{oro} sono le due variabili del candidato $B$ (\texttt{B\_FlightDistance}, \texttt{B\_VertexChi2}); le blu sono per-traccia.
\item Cosa guardare: quali feature dominano in (a) e in (b); se le due graduatorie coincidono.
\end{enumerate}

\sez{Cosa dice}
\begin{enumerate}
\item \textbf{(a) RF.} Dominano gli $\mathtt{IPChi2}$ delle tre tracce: $H2\approx0.35$, $H1\approx0.19$, $H3\approx0.10$ (insieme $\approx60\%$ dell'importanza). Poi $\mathtt{B\_FlightDistance}$ ($\approx0.08$, quarta) e $\mathtt{B\_VertexChi2}$ ($\approx0.03$, quinta). Gli impulsi e i PID pesano ognuno $\lesssim0.03$.
\item \textbf{(b) Lasso.} Vincono $\mathtt{B\_FlightDistance}$ ($\approx0.64$), $\mathtt{H3\_IPChi2}$ ($\approx0.51$), $\mathtt{H2\_PZ}$ e $\mathtt{B\_VertexChi2}$ ($\approx0.36$); anche $\mathtt{H1\_PZ}$ e $\mathtt{H1\_ProbK}$ ($\approx0.29$).
\item \textbf{Accordo parziale.} Spearman $0.42$, $p=0.06$: la RF concentra su $\mathtt{IPChi2}$ (split non lineari), il Lasso distribuisce su più variabili.
\item \textbf{Cosa hanno in comune.} Entrambi mettono le due variabili del candidato $B$ tra le prime cinque: la \emph{topologia di vertice spostato} è la fisica che il modello usa, non la massa.
\item \textbf{Conseguenza.} Il classificatore sfrutta la firma di un vero $B$ (vita media non nulla), come atteso da un discriminante S/B.
\end{enumerate}

\sez{Limiti e orale}
\begin{itemize}
\item L'importanza non è causalità: riflette l'uso nel modello, e con feature correlate la RF la ripartisce arbitrariamente.
\item La prova diretta che il modello non scolpisce la massa è la Fig.~XL (riaddestramento senza impulsi).
\item \textbf{``Perché due misure di importanza?''} Perché rispondono a domande diverse (uso in split non lineari contro effetto marginale lineare); l'accordo parziale è informativo.
\item \textbf{``Cosa fa la penalità $L_1$ che $L_2$ non fa?''} Annulla coefficienti (selezione); $L_2$ li riduce senza annullarli. Qui l'ottimo di $L_1$ non annulla nessun coefficiente (non c'è feature irrilevante).
\end{itemize}

% ======================================================================== XXV

\figh{Fig.\ XXXII --- Scansione della regolarizzazione (App.\ I.6)}{ml_regularization}{0.34}

\sez{Il problema}
La regressione logistica è il modello lineare di riferimento. Quanta regolarizzazione serve, e il confine lineare soffre di varianza (overfitting) o di bias? La scansione di $C$ lo dice.

\sez{Teoria}
Nella regressione logistica regolarizzata si minimizza $\mathcal L(w)+\alpha\,R(w)$, con $R=\|w\|_1$ (Lasso, $L_1$) oppure $\|w\|_2^2$ (Ridge, $L_2$) e $C=1/\alpha$. \emph{Ridge} riduce tutti i coefficienti (shrinkage) senza annullarli; \emph{Lasso} ne azzera alcuni, selezionando le feature. Aumentare $\alpha$ (cioè ridurre $C$) aumenta il bias e riduce la varianza: l'errore di test ha un minimo interno se il modello sovradatta, ed è monotòno (migliore senza penalità) se il problema è il bias.

\sez{Come si legge}
\begin{enumerate}
\item Asse $x$: $C=1/\alpha$ (log, da $10^{-3}$ a $10^{2}$); a destra regolarizzazione debole. Asse $y$: AUC in cross-validation.
\item Due curve: $L_1$ (oro) e $L_2$ (verde). Le linee punteggiate mostrano l'ottimo di ciascuna.
\item Cosa guardare: esiste un massimo interno (segno di overfitting a sinistra)? Le due penalità differiscono, e dove?
\end{enumerate}

\sez{Cosa dice}
\begin{enumerate}
\item \textbf{Nessuna sovra-regolarizzazione necessaria.} Entrambe le curve salgono e si appiattiscono a $\approx0.769$: l'ottimo $L_1$ a $C=0.12$ e $L_2$ a $C=9.2$ danno la stessa AUC ($0.7686$).
\item \textbf{$L_2$ è piatta.} Dal regime più fortemente regolarizzato ($C=10^{-3}$, AUC $\approx0.76$) fino a $C=100$ la differenza è $\sim0.01$: i coefficienti ridotti contano ancora.
\item \textbf{$L_1$ crolla a forte regolarizzazione.} A $C\lesssim3\times10^{-3}$ perde fino a $\sim0.2$ di AUC ($\approx0.55$ a $C=10^{-3}$): i coefficienti vengono azzerati e il modello perde informazione.
\item \textbf{Conseguenza.} L'ottimo $L_1$ non azzera nessun coefficiente: il modello lineare non ha overfitting, il suo limite è il confine lineare (bias), non la varianza. Per questo il salto ai modelli non lineari ($+0.09$ di AUC, Fig.~XXX) è un guadagno di capacità, non di regolarizzazione.
\end{enumerate}

\sez{Limiti e orale}
\begin{itemize}
\item \textbf{``Perché $L_1$ e $L_2$ hanno lo stesso ottimo?''} Perché il problema ha $20$ feature e $\sim10^4$ campioni: la varianza è già piccola, e la regolarizzazione non guadagna nulla.
\item \textbf{``Perché $C$ ottimo è diverso?''} La normalizzazione di $\alpha$ è diversa nelle due penalità; l'importante è l'AUC.
\item \textbf{``Quando serve il Lasso?''} Con molte feature irrilevanti: qui nessun coefficiente è azzerato all'ottimo.
\end{itemize}

\figh{Fig.\ XXXIII --- Bias--varianza del gradient boosting (App.\ I.6)}{ml_nominal_bias_variance}{0.30}
\sez{Il problema}
Quanta complessità è giusta? Un modello troppo semplice sottostima la struttura (bias); uno troppo complesso adatta le fluttuazioni del training (varianza) e generalizza male. Si guarda come cambiano le prestazioni al crescere della complessità, per il modello di riferimento.

\sez{Teoria}
\emph{Decomposizione bias--varianza} dell'errore quadratico atteso a un punto: $\E[(y-\hat f)^2]=\mathrm{Bias}^2[\hat f]+\mathrm{Var}[\hat f]+\sigma^2_{\rm noise}$. Aumentando la complessità il bias scende e la varianza sale; l'errore di test ha tipicamente un minimo.

\emph{Boosting:} $F_m=F_{m-1}+\nu\,h_m$, con $h_m$ un albero (qui profondità $3$) addestrato sul \emph{gradiente negativo} della loss rispetto a $F_{m-1}$; $\nu=0.05$ è il learning rate. Ogni stadio aggiunge complessità: la loss di train scende sempre, quella di test ha un minimo.

\emph{Class weights:} con pesi $w_c=N/(2N_c)$ ($4.40$ per background-like, $0.56$ per signal-like) si minimizza l'errore \emph{pesato}, non il tasso d'errore.

\sez{Come si legge}
\begin{enumerate}
\item Asse $x$ comune: numero di stadi di boosting ($1$--$300$) = complessità del modello.
\item \textbf{(a)} Accuratezza al taglio $0.5$ su train (grigio) e test (oro); la linea punteggiata è ``predici sempre signal-like'' ($0.886$).
\item \textbf{(b)} AUC su train (grigio) e test (oro).
\item Si guarda: \emph{(i)} dove la curva di test smette di salire (punto di saturazione); \emph{(ii)} l'ampiezza del divario train--test (varianza).
\end{enumerate}

\sez{Cosa dice}
\begin{enumerate}
\item \textbf{Regime di bias.} A pochi stadi (fino a $\sim20$--$50$) sia train sia test sono bassi: AUC $\approx0.82$ a $1$ stadio, il modello è troppo semplice.
\item \textbf{Saturazione.} Il test AUC sale fino a $\sim150$ stadi e poi sta a $0.865$; il train continua a salire fino a $\approx0.95$.
\item \textbf{Divario (varianza).} Cresce da $0.018$ a $20$ stadi a $0.084$ a $300$: gli stadi in più adattano le fluttuazioni.
\item \textbf{Nessuna perdita sul test.} L'AUC di test non scende: overfitting benigno, perché il ranking resta stabile; i $200$ stadi di default non lo raggiungono.
\item \textbf{Accuratezza.} Test $\approx0.867$ \emph{sotto} $0.886$: con i pesi il modello minimizza l'errore pesato, in cui un background-like sbagliato costa $\approx8$ volte, non l'accuratezza.
\end{enumerate}

\sez{Limiti e orale}
\begin{itemize}
\item L'accuratezza sotto la baseline non è un errore: la metrica ottimizzata è un'altra; per questo si riporta l'AUC.
\item \textbf{``Come scegli il numero di stadi?''} Si guarda dove il test si satura (o l'early stopping su un set di validazione); qui $\sim150$ e si usa $200$ senza perdita.
\item \textbf{``E per i modelli lineari?''} Il Lasso/Ridge non sovradatta: il limite è il confine lineare (alto bias), non la varianza (l'ottimo $L_1$ non annulla coefficienti).
\item \textbf{``Boosting contro bagging?''} Il boosting riduce il bias (stadi sequenziali sui residui), il bagging (RF) riduce la varianza (alberi indipendenti mediati).
\end{itemize}

% ======================================================================== XXVI
\figh{Fig.\ XXXIV --- Calibrazione (App.\ I.6)}{ml_nominal_calibration}{0.42}
\sez{Il problema}
L'AUC dipende solo dall'\emph{ordinamento} degli score. Se si volesse leggere lo score come una \emph{probabilità} (per pesare i candidati, per esempio) quella probabilità è corretta? E se no, si può correggere senza perdere nulla?

\sez{Teoria}
Un classificatore è \emph{calibrato} se $P(y=1\mid\hat p=s)=s$. \emph{Diagramma di affidabilità:} si raggruppano i candidati per $\hat p$ e si traccia la frazione osservata di positivi. \emph{Brier score:} $\mathrm{BS}=\frac1N\sum_i(p_i-y_i)^2$ (errore quadratico medio rispetto all'etichetta $0/1$, più basso è meglio).

\emph{Effetto dei pesi.} Un classificatore addestrato con pesi $w_s,w_b$ converge a $\tilde p$ tale che
\[
\frac{\tilde p}{1-\tilde p}=\frac{w_s}{w_b}\cdot\frac{p}{1-p}\quad\Longleftrightarrow\quad \tilde p=\frac{w_sp}{w_sp+w_b(1-p)},
\]
cioè le \emph{odds} sono scalate per $w_s/w_b$ ($=0.56/4.40\approx0.13$): il modello descrive un mondo con classi ribilanciate e \emph{sottostima} $p$ del signal-like.

\emph{Calibrazione isotonica:} regressione monotòna non decrescente a gradini (algoritmo PAV) di $y$ su $\hat p$, adattata su fold non usati nel training: corregge non-linearità arbitrarie senza cambiare l'ordinamento.

\sez{Come si legge}
\begin{enumerate}
\item \textbf{Pannello alto:} asse $x$ = probabilità predetta (score grezzo della RF), asse $y$ = frazione osservata di signal-like nel bin. La diagonale punteggiata è la calibrazione perfetta.
\item Due curve: \textbf{grigia} = score grezzo (Brier $0.122$, barre binomiali); \textbf{oro} = ricalibrata isotonicamente (Brier $0.062$).
\item \textbf{Pannello basso:} istogramma (scala log) dello score grezzo per le due classi: signal-like in blu, background-like in arancione.
\item Cosa guardare: \emph{sopra} la diagonale = sottostima, \emph{sotto} = sovrastima; e quanti candidati stanno in ogni bin (istogramma).
\end{enumerate}

\sez{Cosa dice}
\begin{enumerate}
\item \textbf{Score grezzo.} I punti stanno \emph{sopra} la diagonale: a score $0.25$ la frazione osservata è $0.58$, a $0.4$ è $0.75$. Il modello sottostima la probabilità di essere signal-like, nel verso previsto dalla formula dei pesi.
\item \textbf{Ricalibrata.} I punti seguono la diagonale entro le fluttuazioni; il Brier si dimezza da $0.122$ a $0.062$.
\item \textbf{Istogramma.} I signal-like si concentrano a score $0.6$--$0.9$ (fino a $\sim10^3$ per bin), i background-like a score $<0.3$; le code sono popolate da poche decine di candidati.
\item \textbf{Cosa fa la ricalibrazione.} Annulla in gran parte il ribilanciamento dei pesi più che correggere il modello.
\item \textbf{Conseguenza.} Nessun risultato del report cambia: la ricalibrazione è monotòna e il taglio a $0.5$ riguarda lo score grezzo.
\end{enumerate}

\sez{Limiti e orale}
\begin{itemize}
\item I bin a bassa probabilità hanno pochi candidati (barre larghe): la curva è affidabile al centro, meno agli estremi.
\item La ``probabilità'' è quella dell'etichetta debole (finestra di picco contro bande), non di essere un $B$ vero.
\item \textbf{``Perché i class weights?''} Le classi sono sbilanciate ($1940$ contro $15\,130$): senza pesi il classificatore ignorerebbe i background-like.
\item \textbf{``Perché il Brier e non il log-loss?''} Il Brier è limitato e meno sensibile agli score estremi; entrambi sono ``proper scoring rules''.
\end{itemize}

% ======================================================================== XXVII
\figh{Fig.\ XXXV --- Incertezza per candidato (App.\ I.7)}{ml_nominal_uncertainty}{0.28}
\sez{Il problema}
Lo score è un numero. Quanto ci si può fidare di \emph{quel} numero per \emph{quel} candidato? Come si dà un'incertezza a una predizione di un classificatore, e da cosa è fatta?

\sez{Teoria}
Due componenti dell'incertezza predittiva. \emph{Aleatoria:} rumore intrinseco (le classi si sovrappongono nello spazio delle feature): non scende con più dati, è la varianza di Bernoulli $p(1-p)$, massima a $p=0.5$. \emph{Epistemica:} incertezza sul modello dovuta al campione finito: scende con più dati. Per $M$ predizioni $p_m(x)$ la varianza totale si scompone come
\[
\Var[y\mid x]=\underbrace{\E_m[p_m(1-p_m)]}_{\text{aleatoria}}+\underbrace{\Var_m[p_m]}_{\text{epistemica}},
\]
e la dispersione dell'ensemble stima il secondo termine. Gli stimatori:
\begin{center}\small
\begin{tabular}{p{3.2cm}p{6.3cm}p{4.5cm}}
\toprule Stimatore & Come si ottengono le $M$ predizioni & Cosa varia \\ \midrule
Random Forest ($M=300$) & i $300$ alberi, ciascuno su un campione bootstrap & dati e feature \\
MC dropout ($M=300$) & $300$ passaggi della FFNN con dropout \emph{acceso} anche in predizione & struttura della rete \\
Deep ensemble ($M=10$) & $10$ FFNN con inizializzazioni casuali diverse & solo l'inizializzazione \\ \bottomrule
\end{tabular}\end{center}
Il MC dropout (Gal e Ghahramani) è un'approssimazione variazionale di un posterior bayesiano sui pesi: lasciare il dropout attivo equivale a campionare sottoreti. Il deep ensemble (Lakshminarayanan) è un'approssimazione pratica della marginalizzazione sui modi della loss. Una rete bayesiana vera con inferenza variazionale non è stata usata: richiede un'implementazione dedicata.

\sez{Come si legge}
\begin{enumerate}
\item \textbf{(a)} Asse $x$: score medio degli alberi; asse $y$: spread (deviazione standard) tra gli alberi. Curve = mediana per classe (blu signal-like, arancione background-like), banda = $16$--$84\%$.
\item \textbf{(b)} Istogrammi della stima di incertezza per candidato per i tre stimatori (RF verde, MC dropout rosa, deep ensemble marrone); asse $y$: densità.
\item \textbf{(c)} Asse $x$: dimensione del training set ($\approx2500$, $5100$, $10\,200$); asse $y$: spread medio degli alberi, per tutti, background-like e signal-like.
\item Cosa guardare: (a) la \emph{forma} in funzione dello score (campana); (b) le \emph{scale} relative dei tre stimatori; (c) la \emph{pendenza} con la statistica (decrescente = componente epistemica).
\end{enumerate}

\sez{Cosa dice}
\begin{enumerate}
\item \textbf{(a) Aleatoria.} Campana con massimo $\approx0.26$ per score $0.4$--$0.5$ e minimo ($\approx0.06$--$0.09$) agli estremi, uguale per le due classi: quando le classi si confondono gli alberi non concordano; è la forma $p(1-p)$.
\item \textbf{(b) Scale.} Medie $0.178$ (RF, larga $0.05$--$0.35$), $0.095$ (MC dropout; $0.099\pm0.006$ su dieci inizializzazioni), $0.050$ (deep ensemble). Il fattore $\sim3.5$: ciascun albero, di profondità limitata, è uno stimatore povero e discreto; dieci reti sugli stessi dati finiscono su soluzioni simili.
\item \textbf{Correlazioni.} Per candidato $0.41$--$0.57$ (Spearman): gli stessi candidati sono ``difficili'' ma su scale diverse; i metodi non sono intercambiabili.
\item \textbf{(c) Epistemica.} Lo spread RF scende $0.220\to0.204\to0.179$ mentre l'AUC resta $0.860$--$0.862$: più dati riducono l'incertezza sul modello senza cambiare la capacità discriminante.
\item \textbf{Conseguenza.} Nello spread della RF ci sono \emph{entrambe} le componenti; i tre stimatori misurano cose diverse, non va confrontato il valore assoluto.
\end{enumerate}

\sez{Limiti e orale}
\begin{itemize}
\item Non è l'incertezza su $\Araw$ né un intervallo di confidenza sulla probabilità vera: è dispersione d'ensemble, e non entra nei fit.
\item \textbf{``Come dai un'incertezza a una rete neurale?''} MC dropout (approssimazione bayesiana) e deep ensemble (pratico); la RF è la baseline naturale.
\item \textbf{``Perché non una rete bayesiana vera?''} Richiede inferenza variazionale con un'implementazione dedicata; MC dropout dà lo stesso tipo di ensemble a costo nullo.
\item \textbf{``L'incertezza ensemble è una probabilità?''} No, è una deviazione standard tra predizioni; la sua interpretazione probabilistica richiede la teoria dell'inferenza variazionale.
\end{itemize}

% ======================================================================== XXVIII

\figh{Fig.\ XXXVI --- Inerzia e silhouette contro $K$ (App.\ I.8)}{ml_kmeans_selection}{0.30}

\sez{Il problema}
$K$-means richiede $K$ fissato in anticipo. Come si sceglie, e perché nel report si usa $K=8$ mentre il silhouette ha il massimo a $K=3$?

\sez{Teoria}
\emph{Inerzia} $W(K)=\sum_k\sum_{i\in C_k}\|x_i-\bar x_k\|^2$: decresce \emph{sempre} con $K$ (con $K=N$ vale $0$), quindi non ha un minimo: si cerca un ``gomito'' (punto in cui la riduzione rallenta). \emph{Silhouette} $s=\langle(b_i-a_i)/\max(a_i,b_i)\rangle$ misura quanto i cluster sono ben separati ($a_i$ distanza media intra-cluster, $b_i$ dal cluster più vicino); massimizzarla è un criterio di scelta di $K$. Per dati con densità continua e senza cluster naturali (come il piano di Dalitz) entrambi i criteri sono poco decisivi.

\sez{Come si legge}
\begin{enumerate}
\item Asse $x$ comune: $K=2,\dots,12$. \textbf{(a)} Inerzia; \textbf{(b)} silhouette media.
\item Linea tratteggiata rossa: $K=8$ (fissato in anticipo); punto-tratto: $K=3$ (massimo della silhouette).
\item Cosa guardare: (a) c'è un gomito netto? (b) quanto è alto il massimo rispetto agli altri $K$?
\end{enumerate}

\sez{Cosa dice}
\begin{enumerate}
\item \textbf{(a)} L'inerzia scende regolarmente da $\approx2.4\times10^4$ ($K=2$) a $\approx3\times10^3$ ($K=12$) senza gomito netto (leggera flessione intorno a $K=3$--$4$): non indica un $K$ naturale.
\item \textbf{(b)} Silhouette massima a $K=3$ ($0.46$), minimo relativo a $K=4$--$6$ ($\approx0.38$), massimo locale a $K=8$ ($0.41$), poi decresce.
\item \textbf{Scelta.} $K=8$ è un compromesso risoluzione/statistica (almeno $1233$ candidati per cluster), \emph{fissato prima} di guardare le asimmetrie; $K=3$ viene riportato come controllo ($\chi^2=15.2/2$).
\item \textbf{Conseguenza.} Il pattern di asimmetria non dipende dalla scelta: sia $K=3$ sia $K=8$ danno non-uniformità.
\end{enumerate}

\sez{Limiti e orale}
\begin{itemize}
\item La silhouette di $0.46$ non è alta: nessuna struttura a cluster netta; $K$-means è qui una \emph{partizione} adattiva, non una scoperta di gruppi.
\item \textbf{``Perché non $K=3$ come ottimo?''} Poche celle larghe perderebbero risoluzione sulla regione di $\sim1\GeV^2$.
\item \textbf{``Cos'è il gomito?''} Il punto oltre il quale aggiungere cluster riduce poco l'inerzia; qui non c'è.
\end{itemize}

\figh{Fig.\ XXXVII --- $K$-means sul piano di Dalitz (App.\ I.8)}{ml_nominal_kmeans}{0.40}
\sez{Il problema}
La griglia $8\times8$ (Fig.~VII) è arbitraria e i suoi bin hanno statistica molto diversa. Una partizione \emph{adattiva}, determinata dalla densità dei dati, ritrova la stessa struttura? Se sì, il pattern non dipende dalla scelta della griglia.

\sez{Teoria}
\emph{$K$-means} (non supervisionato): partiziona i punti $x_i$ in $K$ cluster $C_k$ minimizzando l'inerzia $W=\sum_k\sum_{i\in C_k}\|x_i-\bar x_k\|^2$. Algoritmo di Lloyd: assegna ogni punto al centroide più vicino, ricalcola i centroidi, ripeti fino alla convergenza (a un minimo \emph{locale}). Le celle sono di Voronoi (confini rettilinei, visibili nella figura); $K$ va fissato in anticipo; le variabili vanno standardizzate.

\emph{Silhouette:} $s_i=(b_i-a_i)/\max(a_i,b_i)$, con $a_i$ distanza media dagli altri punti del cluster e $b_i$ dal cluster più vicino; media in $[-1,1]$. \emph{Adjusted Rand Index} (ARI): accordo tra due partizioni (stabilità); $1$ = identiche.

Poiché il clustering usa solo $(m^2_{\rm lo},m^2_{\rm hi})$ e \emph{non la carica}, i confini non sono scelti guardando l'asimmetria: nessun bias di selezione. L'asimmetria del cluster è $A_k=(N_k^--N_k^+)/N_k$ con $S_k=(A_k-A_{\rm glob})/\sigma(A_k)$ come in Fig.~VII.

\sez{Come si legge}
\begin{enumerate}
\item Assi: $m^2_{\rm lo}$ ($x$, $0$--$15\GeV^2$) e $m^2_{\rm hi}$ ($y$, $0$--$30\GeV^2$) di $K^+K^-$; ogni punto è un candidato.
\item Colori e numeri: ogni cluster ha il colore del suo $S_k$ (barra a destra, saturata a $\pm4$) e il valore scritto nell'etichetta. Blu = deficit di $B^-$ rispetto al globale, rosso = eccesso.
\item I confini sono segmenti rettilinei (celle di Voronoi attorno ai centroidi).
\item Cosa guardare: cluster con $|S_k|>3$ e la loro posizione nel piano.
\end{enumerate}

\sez{Cosa dice}
\begin{enumerate}
\item \textbf{Cluster più significativo.} Blu intenso a basso $m^2_{\rm lo}$ ($1$--$2.5\GeV^2$) e $m^2_{\rm hi}$ intermedio ($\sim7$--$13\GeV^2$): $S=-5.1$, $A_k=-0.127\pm0.017$, $n=3601$.
\item \textbf{Positivi.} Due cluster rossi ad alto $m^2_{\rm hi}$: $S=+3.2$ ($A_k=+0.020\pm0.020$) e $+2.9$ ($+0.040\pm0.028$). Gli altri cinque sono compatibili con $0$ ($|S|\le1.2$).
\item \textbf{Test globale.} $\chi^2=46.4$ per $7$ gradi di libertà, $p=7\times10^{-8}$: non uniforme. Nella finestra di segnale il pattern è più netto ($\chi^2=62.3/7$).
\item \textbf{Robustezza.} $K=3$ (silhouette ottimo $0.46$ contro $0.41$ per $K=8$) dà $\chi^2=15.2/2$; ARI $>0.998$ ripetendo con altra inizializzazione.
\item \textbf{Conseguenza.} Stesso pattern della griglia fissa (negativo a basso $m^2_{\rm lo}$, positivo ad alto $m^2_{\rm hi}$): non è un artefatto della griglia.
\end{enumerate}

\sez{Limiti e orale}
\begin{itemize}
\item I numeri di cluster e griglia non sono indipendenti (stessi dati): l'accordo è sul \emph{pattern}, non sul $p$-value; il $p$-value di $K$-means beneficia di bin più grandi.
\item Stessi caveat: nessuna sottrazione del fondo, nessuna correzione di rivelazione.
\item \textbf{``Perché $K=8$ e non il silhouette ottimo $K=3$?''} Compromesso risoluzione/statistica fissato in anticipo ($\ge1233$ candidati per cluster); se fosse scelto dopo aver visto le asimmetrie ci sarebbe un look-elsewhere. Si riporta anche $K=3$.
\item \textbf{``K-means su dati non gaussiani?''} Assume cluster compatti e simili; qui serve solo come partizione adattiva, non come modello probabilistico.
\end{itemize}

% ======================================================================== XXIX
\figh{Fig.\ XXXVIII --- Soglia sullo score (App.\ I.9)}{ml_nominal_threshold}{0.28}
\sez{Il problema}
Un taglio sul classificatore può modificare l'asimmetria di carica: se la sua efficienza è diversa per $B^+$ e $B^-$ introduce un bias, e l'ML non sarebbe un controllo ma una fonte di artefatto. Lo fa, al variare della soglia?

\sez{Teoria}
Se la selezione ha efficienze $\varepsilon^\pm$ per le due cariche, $A_{\rm sel}=\frac{\varepsilon^-N^--\varepsilon^+N^+}{\varepsilon^-N^-+\varepsilon^+N^+}$, che differisce da $A$ solo se $\varepsilon^-\neq\varepsilon^+$. Il classificatore usa feature che non includono la carica, ma impulsi e PID possono correlarsi con essa: va verificato.

Il campione selezionato è contenuto nel completo (\emph{annidato}), quindi la fluttuazione attesa della differenza è $\sqrt{\sigma^2_{\rm sub}-\sigma^2_{\rm full}}$ (Fig.~XXI), non l'errore del singolo punto. Per la conta, $\sigma(A_{\rm count})=\sqrt{(1-A^2)/N}$.

\sez{Come si legge}
\begin{enumerate}
\item Asse $x$: soglia sullo score del gradient boosting ($0.1$--$0.9$); asse $y$: $\Araw$ di conteggio dei candidati sopra soglia.
\item Punti oro con barre = asimmetria con errore binomiale; linea punteggiata orizzontale = valore della finestra completa ($-0.0432$).
\item Banda grigia = fluttuazione statistica attesa di un sottocampione annidato attorno a quel valore.
\item \textbf{Si confronta ogni punto con la banda}, non con il suo errore: i punti sono correlati (annidati).
\end{enumerate}

\sez{Cosa dice}
\begin{enumerate}
\item \textbf{Stabilità.} I punti stanno tra $\approx-0.040$ e $\approx-0.047$, tutti dentro la banda (entro $1.4\times$ la fluttuazione attesa).
\item \textbf{Allargamento.} Le barre crescono con la soglia: da $18\,443$ a $2621$ candidati; la banda si allarga allo stesso modo.
\item \textbf{Nessuna deriva.} Non c'è una tendenza monotòna: una dipendenza dalla soglia indicherebbe che l'efficienza è carica-dipendente.
\item \textbf{Per polarità} (taglio $0.5$) la conta dà $-0.060\pm0.010$ e $-0.024\pm0.012$ ($2.2\sigma$), ma i fit, che separano il fondo, $-0.054$ e $-0.032$ ($1.3\sigma$): la maggior parte della differenza di conteggio viene dal fondo.
\item \textbf{Conseguenza.} La selezione non introduce un bias rivelabile nell'asimmetria di conteggio.
\end{enumerate}

\sez{Limiti e orale}
\begin{itemize}
\item È il conteggio, che include il fondo: la verifica sul fit è nella Fig.~XXXIX.
\item \textbf{``Perché non l'errore del singolo punto?''} Perché i campioni sono contenuti l'uno nell'altro: la differenza ha varianza $\sigma^2_{\rm sub}-\sigma^2_{\rm full}$, non la somma.
\item \textbf{``Cosa vorrebbe dire una deriva?''} Un'efficienza carica-dipendente: il classificatore userebbe informazione correlata con la carica (impulsi, per via dell'asimmetria di rivelazione).
\item Con $9$ soglie correlate, un punto a $1.4\times$ la banda non è significativo.
\end{itemize}

% ======================================================================== XXX
\figh{Fig.\ XXXIX --- Tre metodi, campione completo contro selezionato (App.\ I.9)}{ml_nominal_three_methods}{0.30}
\sez{Il problema}
Il taglio ML cambia il risultato \emph{con tutti e tre} i modi di stimare $\Araw$, ciascuno con i suoi assunti? Se uno solo cambiasse, il segnale sarebbe un artefatto del metodo.

\sez{Teoria}
Tre stimatori: (i)~fit bayesiano (mediana a posteriori); (ii)~fit frequentista (massimo di verosimiglianza); (iii)~conteggio di tutti i candidati. Il fit stima l'asimmetria del \emph{segnale} separando il fondo; il conteggio misura $A_{\rm count}=f_sA+(1-f_s)A_b$ (con $f_s$ la frazione di segnale): include il fondo, ha varianza minore perché non ha incertezze di modello, ma è polarizzato da $A_b$. Lo spostamento atteso tra campione completo e annidato è $\sqrt{\sigma^2_{\rm sub}-\sigma^2_{\rm full}}$.

\sez{Come si legge}
\begin{enumerate}
\item Asse $x$: $\Araw$ (da $-0.052$ a $-0.036$); tre righe, una per metodo (Bayesiano in alto, frequentista, conteggio).
\item In ogni riga: punto nero con barra = campione completo ($19\,731$); punto oro = campione con score $>0.5$ ($15\,709$).
\item La banda grigia verticale è $\pm1\sigma$ del bayesiano sul campione completo; la linea punteggiata il valore centrale.
\item Cosa guardare: il punto oro cade dentro la banda? Le barre oro e nere si sovrappongono?
\end{enumerate}

\sez{Cosa dice}
\begin{enumerate}
\item \textbf{Bayes:} $-0.0443\pm0.0085\to-0.0451\pm0.0087$ (spostamento $-0.0008$; atteso $0.0018$).
\item \textbf{Frequentista:} $-0.0441\pm0.0085\to-0.0447\pm0.0086$ ($-0.0006$; atteso $0.0012$).
\item \textbf{Conteggio:} $-0.0432\pm0.0071\to-0.0446\pm0.0080$ ($-0.0014$; atteso $0.0036$).
\item \textbf{Significatività.} $5.2\sigma$ in entrambi i campioni (profilo di verosimiglianza).
\item \textbf{Precisione.} Invariata ($0.0087$ contro $0.0085$): il fondo sotto il picco è già piccolo ($\approx4\%$ nella finestra $\pm2\sigma$); il taglio tiene il $92\%$ del segnale ($N_s=14\,121$) e il $36\%$ del fondo ($N_b=1590$).
\item \textbf{Conseguenza.} Tre metodi con assunzioni diverse danno lo stesso risultato sul campione ML-selezionato: $\Araw$ non è un artefatto della selezione o del metodo.
\end{enumerate}

\sez{Limiti e orale}
\begin{itemize}
\item Il conteggio ha barra più corta ma non è comparabile: include il fondo, il suo valore non è l'asimmetria del segnale.
\item \textbf{``Perché la precisione non migliora?''} Perché c'è poco fondo da togliere; il selezionato ha statistica minore ($15\,709$ contro $19\,731$) e quindi errore leggermente maggiore.
\item \textbf{``Cosa avrebbe significato un'asimmetria diversa?''} Un'efficienza di selezione carica-dipendente, o un fondo con asimmetria diversa in quel campione.
\item Gli errori dei due campioni sono correlati (annidati): non si sommano in quadratura.
\end{itemize}

% ======================================================================== XXXI
\figh{Fig.\ XL --- Massa dopo il taglio ML (App.\ I.9)}{ml_nominal_mass_after_cut}{0.28}
\sez{Il problema}
Il conteggio e i fit non dicono se il classificatore \emph{deforma lo spettro di massa}. Lo deforma, e perché? Se il modello avesse imparato la massa dagli impulsi, un taglio sullo score creerebbe un picco finto (``sculpting'') e la selezione sarebbe invalida per qualunque misura basata sulla massa.

\sez{Teoria}
\emph{Sculpting}: se lo score è correlato con la massa, un taglio modella la distribuzione di massa del fondo (può creare un picco finto). Condizione di CWoLa: le feature non devono contenere informazione sulla massa. Ma la massa è funzione delle \emph{nove componenti di impulso} delle tre tracce ($m^2=(\sum E)^2-|\sum\vec p|^2$), quindi un modello flessibile potrebbe impararla: il test è riaddestrare \emph{senza} le componenti di impulso e confrontare.

\emph{Efficienza di taglio} in funzione della massa: $\varepsilon(m)=n_{\rm sel}(m)/n_{\rm all}(m)$.

\sez{Come si legge}
\begin{enumerate}
\item Asse $x$ comune: $m(KKK)$ ($5100$--$5600\MeV$).
\item \textbf{(a)} $y$ = candidati per $25\MeV$ in scala log: nero = tutti; oro pieno = score $>0.5$ del classificatore a $20$ feature; oro tratteggiato = score $>0.5$ del classificatore \emph{senza} impulsi.
\item \textbf{(b)} $y$ = $\varepsilon(m)$: frazione che passa il taglio (stessi due classificatori).
\item Zone: a righe = $5100$--$5150\MeV$ (esclusa dalle etichette di fondo); grigio = finestra signal-like ($\pm2\sigma$).
\item Cosa guardare: (i) le due curve oro coincidono? (ii) $\varepsilon(m)$ è piatta o ha un picco in corrispondenza del segnale?
\end{enumerate}

\sez{Cosa dice}
\begin{enumerate}
\item \textbf{No sculpting.} Con e senza impulsi le curve coincidono (AUC $0.864$ contro $0.860$): il classificatore \emph{non} ricostruisce la massa dagli impulsi.
\item \textbf{Efficienza.} $\varepsilon\approx0.89$--$0.90$ al picco, $\approx0.70$ sotto $5150\MeV$, $\approx0.3$--$0.4$ a $5150$--$5250\MeV$ e $0.06$--$0.13$ sopra $5375\MeV$.
\item \textbf{Cosa seleziona.} Il taglio elimina il combinatorio (massa alta) e tiene i $B$ veri, anche parzialmente ricostruiti (massa bassa, vertice spostato): la \emph{composizione} del fondo cambia.
\item \textbf{Effetto sul modello.} Il $\chi^2$ del modello nominale peggiora ($164\to253$ bayesiano, $442\to557$ frequentista); la pendenza dell'esponenziale lenta passa da $0.0028$ a $0.0071\MeV^{-1}$ ($B^+$). Una Crystal Ball dà lo stesso peggioramento ($407\to522$): non sono le code del segnale.
\item \textbf{Conseguenza.} $\Araw$ non è affetto (picco e fondo libero nel fit), ma il campione selezionato è descritto \emph{peggio} dal modello nominale.
\end{enumerate}

\sez{Limiti e orale}
\begin{itemize}
\item Un'analisi che si basasse sulla forma del fondo, o su un fondo preso dalle bande del campione completo, sarebbe sbagliata dopo il taglio.
\item \textbf{``Come sai che non è sculpting?''} Perché il classificatore senza impulsi (che non può calcolare la massa) dà la stessa efficienza e lo stesso $\chi^2$.
\item \textbf{``Perché l'efficienza è alta sotto $5150\MeV$ se sono fondo?''} Sono in gran parte $B$ parzialmente ricostruiti: hanno vertice spostato come il segnale, e il classificatore non li distingue; le etichette del fondo escludono proprio questa regione.
\item \textbf{``Questo invalida l'ML?''} No per $\Araw$, sì per qualunque misura che usi il fondo.
\end{itemize}

% ======================================================================== XXXII
\figh{Fig.\ XLI --- Campione ad alto fondo (App.\ I.9)}{ml_nominal_high_background}{0.30}
\sez{Il problema}
Nel campione nominale c'è poco fondo, quindi il ML non può migliorare. Dove il fondo è maggiore (identificazione più larga), il classificatore migliora la precisione e rimuove il bias dal fondo mal identificato meglio del taglio sul PID?

\sez{Teoria}
\emph{Fondo piccato} (peaking background): decadimenti mal identificati, per esempio $B^\pm\to\pi^\pm K^+K^-$ con un pione assegnato come kaone, hanno massa ricostruita \emph{vicina} a $m_B$ e stanno sotto il picco; la loro asimmetria di carica è diversa e sposta $\Araw$. Un classificatore CWoLa è addestrato a separare picco da bande: per costruzione \emph{non} può rigettare un fondo che sta sotto il picco, perché le etichette lo marcano come segnale. La particle identification (PID, $\mathtt{ProbK}$) lo rigetta usando informazione indipendente dalla massa.

\emph{Efficienza di segnale} qui = $N_s$ fittato relativo a $\mathtt{ProbK}>0.3$.

\sez{Come si legge}
\begin{enumerate}
\item Campione con $\mathtt{ProbK}>0.3$ ($33\,328$ candidati, purezza $94\%$ nella finestra $\pm2\sigma$). Asse $x$ comune: efficienza di segnale, da $\approx0.15$ a $1$.
\item Due modi di ``pulire'': \textbf{blu} = stringere $\mathtt{ProbK}$ ($0.3,\dots,0.8$); \textbf{oro} = tagliare lo score di un classificatore ririaddestrato su questo campione ($0,\dots,0.9$, AUC $0.895$).
\item \textbf{(a)} $y=\sigma(\Araw)$ (incertezza statistica). \textbf{(b)} $y=\Araw$ del fit frequentista, con barre.
\item La linea punteggiata verticale è la selezione nominale ($\mathtt{ProbK}>0.5$, efficienza $\approx0.68$).
\item Cosa guardare: (a) quale curva sta più in basso a pari efficienza; (b) a quale valore va $\Araw$ al ridursi del fondo mal identificato.
\end{enumerate}

\sez{Cosa dice}
\begin{enumerate}
\item \textbf{(a) Precisione.} A pari efficienza il classificatore dà $\sigma$ leggermente minore: $0.0072\to0.0070$ al meglio (cut a $0.5$: $93\%$ del segnale, purezza $99\%$). Il guadagno è piccolo.
\item \textbf{(b) Valore.} Il campione largo dà $\Araw\approx-0.055$; stringere $\mathtt{ProbK}$ lo riporta a $\approx-0.047$ (a $0.4$) e $\approx-0.044$ (nominale): spostamento $+0.008$, $3.3$ volte lo spread atteso.
\item \textbf{Il taglio sullo score non lo sposta.} Resta a $\approx-0.055$ (spostamento $-0.001$) anche con il $93\%$ del segnale.
\item \textbf{Perché.} Il fondo mal identificato (peaking) ha massa vicina al $B$: le etichette basate sulla massa lo marcano come signal-like e il classificatore non lo può rigettare.
\item \textbf{Conseguenza.} Il ML migliora poco la precisione e non rimuove il fondo che conta per $\Araw$; la particle identification resta necessaria, e il valore nominale ($\mathtt{ProbK}>0.5$) è giustificato.
\end{enumerate}

\sez{Limiti e orale}
\begin{itemize}
\item Errori grandi e campioni annidati (punti correlati): la tendenza è chiara, ma i singoli punti no.
\item \textbf{``Come lo si controllerebbe?''} Servirebbero etichette con informazione oltre la massa (simulazione) o un modello del fondo piccato.
\item \textbf{``Perché l'efficienza di segnale è relativa?''} Perché non c'è un campione di riferimento a efficienza nota; si usa $N_s$ fittato.
\item \textbf{``Differenza tra PID e classificatore?''} La PID usa informazione di rivelatori dedicati (RICH, calorimetri) indipendente dalla massa; il classificatore usa cinematica e vertice, che per costruzione non distingue il peaking.
\end{itemize}


\figh{Fig.\ XLII --- Ricerca dell'architettura della FFNN (App.\ I.10)}{ml_ffnn_search}{0.34}

\sez{Il problema}
La rete ha iperparametri (architettura, dropout, learning rate). I default sono vicini all'ottimo? E quanto conta l'architettura rispetto agli altri?

\sez{Teoria}
La ricerca è una \emph{griglia} di $168$ configurazioni ($7$ architetture $\times$ $4$ dropout $\times$ $6$ learning rate), valutate con cross-validation ripetuta. Il massimo di molte stime rumorose è \emph{polarizzato verso l'alto} (winner's curse): l'AUC di CV della configurazione migliore sovrastima la sua prestazione, per questo la scelta si conferma su un test non usato. Quando molte configurazioni sono nella stessa fascia statistica l'ottimo è un \emph{plateau}.

\sez{Come si legge}
\begin{enumerate}
\item Asse $x$: le sette architetture, dalla più piccola (un solo strato da $8$ unità) alla più grande ($[128,64,32,16,8]$).
\item Ogni punto è una configurazione (dropout e learning rate diversi), con un piccolo spostamento orizzontale per leggerli; il segmento nero è il massimo di ciascuna architettura.
\item Cosa guardare: come cresce il massimo con la capacità? Quanto è larga la dispersione dentro ogni architettura?
\end{enumerate}

\sez{Cosa dice}
\begin{enumerate}
\item \textbf{Capacità.} Il massimo sale da $0.839$ (un strato da $8$ unità) a $0.847$ ($1$ strato più grande), $0.856$ ($2$ strati), $0.858$ (due strati larghi e tre strati) e poi si assesta: $0.857$ (4 strati), $0.854$ (5): non serve più capacità.
\item \textbf{Dispersione.} All'interno di ogni architettura i punti si estendono su $0.02$--$0.09$: la scelta di learning rate e dropout conta più dell'architettura oltre due strati.
\item \textbf{Plateau.} Le prime cinque configurazioni stanno entro $0.0007$ in AUC (molto meno della dispersione tra i fold, $0.017$): l'ottimo è un plateau.
\item \textbf{Conseguenza.} I default sono vicini all'ottimo: il guadagno sul test è $0.004$, sotto l'incertezza ($0.012$); si usano i default.
\end{enumerate}

\sez{Limiti e orale}
\begin{itemize}
\item La ricerca sulla RF ($120$ configurazioni) e sul gradient boosting ($100$) è nei numeri delle tabelle, senza una figura analoga perché i punti per configurazione non sono salvati.
\item \textbf{``Perché non una ricerca bayesiana?''} Con griglie di $\sim100$ configurazioni e $\sim30\,$s per fit è fattibile esaustiva.
\item \textbf{``Il test è stato toccato?''} No: la scelta è in CV e il test confronta il vincitore con il default.
\end{itemize}

\clearpage
\section*{Incongruenze tra figure e testo da sistemare}
\begin{enumerate}
\item \textbf{Fig.\ II} (Marco): il titolo del grafico dice ``toys'', il testo dice repliche bootstrap. Già nella nota per Marco.
\item \textbf{Fig.\ XVIII e XXII}: legende e assi dicono ``pseudo-datasets'', il report usa ``toy MC''.
\item \textbf{Fig.\ XXI}: la normale è centrata a zero mentre l'istogramma è centrato sullo spostamento dei dati; la didascalia dovrebbe dire che si guarda la larghezza.
\item \textbf{Fig.\ XXIII(c)} ripete la Fig.\ VI: è voluto (la VI è nel corpo), ma va dichiarato.
\item \textbf{App.\ C.3} (testo): afferma accordo ``entro il $10\%$ per tutti i parametri di forma; per $\lambda_1^+$ il bootstrap differisce del $26\%$ (Fig.\ XII).
\end{enumerate}
\end{document}
